\documentclass[draft]{amsart}

\usepackage{pgf,tikz}
\usetikzlibrary{arrows}
\usepackage{cite}
\usepackage{blindtext}%\usepackage{eufrak}
\usepackage{amssymb,graphicx,enumerate,color}
\usepackage{multirow}
\usepackage{multicol}%\usepackage{comment}
\usepackage{xcolor}
\renewcommand{\subjclassname}{%
  \textup{2013} Mathematics Subject Classification}


\setlength{\textheight}{220mm} \setlength{\textwidth}{155mm}
\setlength{\oddsidemargin}{1.25mm}
\setlength{\evensidemargin}{1.25mm} \setlength{\topmargin}{0mm}


\renewcommand{\P}{\mathcal{P}}
\newcommand{\U}{\mathcal{U}}
\newcommand{\D}{{\mathcal{D}}}
\renewcommand{\O}{{\mathcal{O}}}
\newcommand{\R}{{\mathcal{R}}}
\newcommand{\M}{\mathcal{M}}
\newcommand{\G}{\mathcal{G}}
\renewcommand{\t}{\varphi}
\newcommand{\x}{\xi}
\newcommand{\z}{\zeta}
%\newcommand{\CCal}{{\mathcal{C}}}
\newcommand{\E}{{\mathcal{E}}}
\newcommand{\T}{{\mathcal{T}}}
\newcommand{\JCal}{{\mathcal{J}}}
\newcommand{\ZBbb}{{\mathbf{Z}}}
\newcommand{\CBbb}{{\mathbf{C}}}
\newcommand{\K}{{\mathcal{K}}}
\newcommand{\qu}{{\underline{q}}}
\newcommand{\thre}{{\sqrt{3}}}
\newcommand{\qq}{{\(\dfrac q{\qu}\)^2}}
\renewcommand{\i}{\infty}
\renewcommand{\b}{\beta}
\def\receivedline{\relax}
\def\dedication#1{\receivedline\vskip4pt
{\normalsize\begin{center}#1\end{center}\vskip1sp}}

\newcommand{\beqs}{\begin{equation*}}
\newcommand{\eeqs}{\end{equation*}}
\numberwithin{equation}{section}
 \theoremstyle{plain}
\newtheorem{theorem}{Theorem}[section]
\newtheorem{lemma}[theorem]{Lemma}
\newtheorem{corollary}[theorem]{Corollary}
\newtheorem{proposition}[theorem]{Proposition}
%\newtheorem{remarks}[theorem]{Remarks}
\newtheorem{entry}[theorem]{Entry}
\theoremstyle{remark}
\newtheorem{remark}{Remark}
\newtheorem{definition}[theorem]{Definition}
\newtheorem{conjecture}[theorem]{Conjecture}
\newtheorem{example}[theorem]{Example}
\newcommand{\divides}{\mid}


\hyphenation{auto-maton}
%\newenvironment{pf}
%   {\vskip 0.15in \par\noindent{\it Proof of}\hskip 0.5em\ignorespaces}
%   {\hfill $\Box$\par\medskip}

\allowdisplaybreaks

\begin{document}

\makeatletter
\def\imod#1{\allowbreak\mkern10mu({\operator@font mod}\,\,#1)}
\makeatother

\author{Xinhua Xiong}
   \address{Research Institute for Symbolic Computation, Johannes Kepler University, A4040 Linz, Austria}
   \email{xxiong@risc.uni-linz.ac.at}




\title[ Euler's partition theorem]{\hspace{0.0cm}  Euler's partition theorem for all moduli and new companions to Rogers-Ramanujan-Andrews-Gordon identities}

\thanks{The  author was  supported  by the Austria Science Foundation (FWF) grant SFB F050-06 (Special Research Programm `` Algorithmic and Enumerative Combinatorics'').}    

\dedicatory{}     

\begin{abstract}  
In this paper, we give a conjecture, which generalises Euler's partition theorem involving odd parts and different parts for all moduli.  We prove this conjecture for two family partitions. We give $q$-difference equations for the related generating function if the moduli is three. We provide new companions to Rogers-Ramanujan-Andrews-Gordon identities under this conjecture.
\end{abstract}   
   
\keywords{Euler's partition theorem, Rogers-Ramanujan-Andrews-Gordon identities, partition identities, generating functions, q-series}
 
 
 \subjclass[2010]{05A17, 11P84}
 %%%%	05A15  	Exact enumeration problems, generating functions [See also 33Cxx, 33Dxx]
 %	05A17  	Partitions of integers [See also 11P81, 11P82, 11P83]
 %  05A19  	Combinatorial identities, bijective combinatorics
%	11B75  	Other combinatorial number theory
 %	11P81  	Elementary theory of partitions [See also 05A17]
 %	11P84  	Partition identities; identities of Rogers-Ramanujan type
 %	33D15  	Basic hypergeometric functions in one variable, ${}_r\phi_s$
 

%\date{June 1, 2016}
   
\maketitle
\section{Introduction}
In the theory of partitions, Euler's partition theorem involving odd parts and different parts is one of the famous theorems. It claims that the number of partitions of $n$ into odd parts is equal to  the number of partitions $n$ into different 
parts. By constructing a bijection, Sylvester \cite{Syl82} not only proved Euler's theorem, but also provided a refinement of it which can be stated as the number of partitions of $n$ into odd parts with exactly $k$ different parts is equal to the number of partitions of $n$ into different parts such that exactly $k$ sequences of consecutive integers occur in each partition.  Bessenrodt  \cite{bes94} proved that Sylvester's bijection implies that the number of partitions of $n$ into different parts with the alternating sum $\Sigma$ is equal to the number of partitions of $n$ with $\Sigma$ odd parts. Here for a partition $\lambda=(\lambda_1,\lambda_2,\dots)$, the alternating sum is defined by 
$$
\Sigma=\lambda_{1}-\lambda_{2}+\lambda_{3}-\lambda_{4}+\dots.
$$
From the theorems on lecture hall partitions, Bousquet-M\'{e}lou and Ericsson \cite{bou97, bou972} also obtained this result. In \cite{Kim-Yee}, Kim and Yee gave a different description of the Sylvester's bijection which provide a more simple proof of the refinement of Euler's theorem due to Bessenrodt, Bousquet-Melou and Ericsson. There are several other refinements and variants of Euler's theorem. See \cite{All3, AndrewsEuler, Chen, Pak, SavageYee,Yee2001, Yee2002, Zeng}.

We can think Euler's theorem as a theorem on partitions about moduli two by interpreting odd parts as  parts $\equiv 1\pmod 2.$ 
The first nontrivial generalisation of Euler's theorem to all moduli in this sense is the following theorem due to Pak-Postinikov.

 \begin{theorem}[Pak-Postinikov \cite{Pak-Postnikov}]   \label{puretype thm} 
 The number of partitions of $n$ with type $(c, m-c, c, m-c,\dots)$  is equal to the number of partitions of $n$ with all parts $\equiv c \pmod m$.
 \end{theorem}
By the type $(c, m-c,c,m-c,\dots)$ for a partition $\lambda$, it means that $\lambda$ has the length divisible by $m$ by allowing zero as parts  and has $c\geq 1$ largest part, $m-c$ second large part, etc. So it has the form:
\begin{align*}
\lambda_{1}=\lambda_{2}=\lambda_{3}=\dots=\lambda_{c}>\lambda_{c+1}=\lambda_{c+2}= \dots = \lambda_{m}>\lambda_{m+1}=\dots=\lambda_{m+c}>\dots
\end{align*}

In this paper, we give a conjecture, which generalises Euler's partition theorem mentioned above to all moduli. Also it generalises 
Pak-Postinikov's theorem. We prove this conjecture for two family partitions.  It is interesting that this conjecture provides new companions to Rogers-Ramanujan-Andrews-Gordon identities. 
 
In Section 2, we state our conjecture and give examples to illustrate our conjecture. In Section 3, we prove the conjecture for two family partitions. In Section 4, we use $q$-difference equations method to study the conjecture for moduli three. In Section 5, we provide new companions to Rogers-Ramanujan-Andrews-Gordon identities under the conjecture. Last section contains discussions of combinatorial arguments behind our proofs. 



\section{Main conjecture}
In order to state our conjecture, we introduce some notations and terminologies.

$\bullet$  For a partition $\lambda=(\lambda_1,\lambda_2, \lambda_{3},\dots,\lambda_{r})$ of $n$, we denote $n$ by $|\lambda |.$ We sometimes write a partition as the form
$$\lambda_1+\lambda_2+\lambda_3+\dots+\lambda_{r}  \quad \mbox{\quad or the form\quad} \quad   \lambda_1\geq\lambda_2\geq\lambda_{3}\geq\dots\geq\lambda_{r}.$$

$\bullet$ Let $m\geq 2$ be an integer, which is the moduli in our sense.  For a partition $\lambda_1\geq\lambda_2\geq\lambda_{3}\geq\dots\geq \lambda_{km}$ ($k\geq1$ ) with length divisible by $m$, we define its alternating sum type to be a $(m-1)$-tube non-negative integers sequence  $(\Sigma_{1},  \Sigma_{2},  \dots, \Sigma_{m-2},  \Sigma_{m-1})$ by
$$
\Sigma_{1}=\sum_{i=1}^{k}\lambda_{(i-1)m+1}-\lambda_{(i-1)m+2}, \Sigma_{2}=\sum_{i=1}^{k}\lambda_{(i-1)m+2}-\lambda_{(i-1)m+3}, \dots, \Sigma_{m-1}=\sum_{i=1}^{k}\lambda_{(i-1)m+m-1}-\lambda_{im}.
$$
For example, if $m=3$, the partition $6+5+4+3+2+1$ has the alternating sum type $(6-5+3-2,5-4+2-1)=(2, 2)$. For the partitions appearing in the theorem of Pak-Postnikov, their alternating sum types are special $(m-1)$-tube integers : $(0,0,\dots,0, \Sigma, 0,\dots,0)$ and $\Sigma$ lies at the $c^{th}$ position.


$\bullet$ For a partition $\lambda_1\geq\lambda_2\geq\lambda_{3}\geq\dots\geq \lambda_{km}$ ($k\geq1$ ),
let $1 \leq i\leq k$, then a partition piece $\lambda_{(i-1)m+1}\geq \lambda_{(i-1)m+2}\geq \dots\geq \lambda_{(i-1)m+m-1}\geq \lambda_{im}$ is called a basic unit of $\lambda$.

$\bullet$ When we speak of alternating sum types, we allow last few parts to be zero so that any partition has length divisible by $m$. For example if $m=3$, the partition $5+4+3+3$ has length $6$ by viewing it as $5+4+3+3+0+0$ and it has two basic units: $5+4+3$ and $3+0+0.$

$\bullet$ For a partition $\lambda$ with all parts $\not\equiv 0\pmod m$, we define its length type to be a $(m-1)$-tube non-negative
integer sequence $(l_{1}, l_{2}, l_{3}, \dots, l_{m-2},  l_{m-1})$, where for $1\leq i\leq m-1$,  $l_{i}$ is the number of parts
of $\lambda$ which are congruent to $i$ modulo $m$. For example, for a partition $\lambda$ with all parts are congruent to $c$, its length type is $(0,0,\dots,0,l,0\dots,0)$, where $l$ is the number of parts of $\lambda$ and lies at the $c^{th}$ position.

$\bullet$ Given a alternating sum type $(\Sigma_{1},\Sigma_{2}, \dots,\Sigma_{m-2}, \Sigma_{m-1})$, if only one $\Sigma_{i}$ is non-zero, we call it a pure alternating sum type, simply a pure type, otherwise, we call it a mixed alternating sum type, simply a mixed type. We have  similar notions of pure types and mixed types for length types.

Now we state our conjecture.
\begin{conjecture}
Let $m\geq 2$,  let $P$ be the set of  partitions with each part can be repeated at most $m-1$ times, this implies their alternating sum types can not be $(0,0,\dots,0)$. Let $Q$ be the set of partitions with no parts $\equiv 0\pmod m$. Then we have
the partition identity:
$$
\sum_{\lambda \in P} z_{1}^{\Sigma_{1}(\lambda)}  z_{2}^{\Sigma_{2}(\lambda)      }  \dots z_{m-1} ^{\Sigma_{m-1}(\lambda)          }  q^{|\lambda |}      =        \sum_{\mu \in Q} z_{1}^{l_{1}(\mu)}  z_{2}^{l_{2}(\mu)      }  \dots z_{m-1} ^{l_{m-1}(\mu)          }  q^{|\mu |}.   
 $$
Equivalently,  the number of partitions  of $n$ with the alternating sum type $(\Sigma_{1}, \Sigma_{2}, \dots \Sigma_{m-1})$
is equal to the number of partitions of $n$ with $\Sigma_{1}$ parts congruent to $1$ modulo $m$, $\Sigma_{2}$ parts congruent to $2$ modulo $m$, $\dots$ , $\Sigma_{m-1}$ parts congruent $m-1$ modulo $m$.
\end{conjecture}

 If we let $z_{1}=z_{2}=\dots z_{m-1}=z$, we get the result that the number of partitions of $n$ with parts repeated at most
 $m-1$ times and total alternating sum $\Sigma_{1}+\Sigma_{2}+\dots +\Sigma_{m-1}$ is equal to the number of partitions
 of $n$ with no parts congruent to $0$ modulo $m$  and $\Sigma_{1}+\Sigma_{2}+\dots +\Sigma_{m-1}$ parts, which  is the refinement of a Glaisher's theorem:
\begin{theorem}[Glaisher \cite{Pak}]
The number of partitions of $n$ with parts repeated  at most $m-1$ times is equal to the number of partitions of $n$ with no
parts is congruent to $0$ modulo $m$.
\end{theorem}
When the alternating sum type is pure type, this conjecture is reduce to Theorem 1.1 due to Pak-Postinikov.
When $m$ is $2$, this conjecture is reduce to the refinement of Euler's theorem due to Bessenrodt, Kim and Yee.

We give $n=10$ and $m=3, 4$ to illustrate this conjecture. 
We list partitions in $P$, their alternating sum types and the numbers on the left, and the corresponding parts for partitions in $Q$ on the right. We only list all partitions with mixed types.
%%%%%%%%%%%%%%%%%%%%%%%%%%%%%%%%%%%%%%%%
%%%%%%%%%%%%%%%%%%%%%%%%%%%%%%%%%%%%%%%%
%\begin{example}$m=3$, $n=11$. 

\begin{minipage}{0.38\textwidth}
\small
\begin{align*}
& \mbox{ Partitions in $P$} &(\Sigma_{1}, \Sigma_{2})&\quad\sharp \\
&\left\{
\begin{array}{llll}
&3+3+2+2+1\\
&5+4+2\\
&4+4+2+1\\
&4+3+2+1+1
\end{array}
\right\} &(1, 2)   &\quad4 \\
&\left\{
\begin{array}{llll}
&6+3+2\\
&5+3+2+1\\
&5+2+2+1+1\\
&4+3+2+2
\end{array}
\right\} &(3,1)  &\quad4\\
&\left\{
\begin{array}{ll}
&6+4+1\\
&5+4+1+1
\end{array}
\right\} &(2,3)  &\quad2 \\
&\left\{
\begin{array}{ll}
&7+3+1\\
&6+3+1+1
\end{array}
\right\} &(4,2)  &\quad2\\
&\left\{
\begin{array}{ll}
&8+2+1\\
&7+2+1+1
\end{array}
\right\} &(6,1)  &\quad2\\
&\left\{
\begin{array}{l}
6+5
\end{array}
\right\} &(1,5)  &\quad1\\
&\left\{
\begin{array}{l}
7+4
\end{array}
\right\} &(3,4)  &\quad1\\
&\left\{
\begin{array}{l}
8+3
\end{array}
\right\} &(5,3)  &\quad1\\
&\left\{
\begin{array}{l}
9+2
\end{array}
\right\} &(7,2)  &\quad1\\
&\left\{
\begin{array}{l}
10+1
\end{array}
\right\} &(9,1)  &\quad1
 \end{align*}
\end{minipage}
%%
\begin{minipage}{0.50\textwidth}
\small
\begin{align*}
& \mbox{ Partitions in $Q$ } &(l_{1}, l_{2})  &\quad\sharp\\
&\left\{
\begin{array}{llll}
&8+2+1\\
&7+2+2\\
&5+5+1\\
&5+4+2
\end{array}
\right\} &(1,2)  &\quad4\\
&\left\{
\begin{array}{llll}
&8+1+1+1\\
&7+2+1+1\\
&5+4+1+1\\
&4+4+2+1
\end{array}
\right\} &(3,1)  &\quad4\\
&\left\{
\begin{array}{ll}
&5+2+2+1+1\\
&4+2+2+2+1
\end{array}
\right\} &(2,3)  &\quad2\\
&\left\{
\begin{array}{ll}
&5+2+1+1+1+1\\
&4+2+2+1+1+1
\end{array}
\right\} &(4,2)  &\quad2\\
&\left\{
\begin{array}{ll}
&5+1+1+1+1+1+1\\
&4+2+1+1+1+1+1
\end{array}
\right\} &(6,1)  &\quad2\\
&\left\{
\begin{array}{l}
2+2+2+2+2+1
\end{array}
\right\} &(1,5)  &\quad1\\
&\left\{
\begin{array}{l}
2+2+2+2+1+1+1
\end{array}
\right\} &(3,4)  &\quad1\\
&\left\{
\begin{array}{l}
2+2+2+1+1+1+1+1
\end{array}
\right\} &(5,3)  &\quad1\\
&\left\{
\begin{array}{l}
2+2+1+1+1+1+1+1+1
\end{array}
\right\} &(7,2)  &\quad1\\
&\left\{
\begin{array}{l}
2+1+1+1+1+1+1+1+1+1
\end{array}
\right\}&(9,1)  &\quad1
\end{align*}  
\end{minipage}
%\end{example}



%%%%%%%%%%%%%%%%%%%%%%%%%%%%%%%

%\vspace{0.5cm}

%\begin{example}
%$m=4$, $n=10$. 

\begin{minipage}{0.43\textwidth}
\small
\begin{align*}
\small
& \mbox{ Partitions in $P$} &(\Sigma_{1}, \Sigma_{2},\Sigma_{3})  &\sharp \\
&\left\{
\begin{array}{lll}
&4+3+2+1\\
&3+3+2+1+1\\
&3+2+2+1+1+1
\end{array}
\right\}\, &(1,1,1) \,&3\\
&\left\{
\begin{array}{ll}
&5+3+1+1\\
&4+3+1+1+1
\end{array}
\right\}\, &(2,2,0) \, &2\\
&\left\{
\begin{array}{ll}
&5+2+2+1\\
&4+2+2+1+1
\end{array}
\right\}\, &(3,0,1) \, &2\\
&\left\{
\begin{array}{ll}
&6+2+1+1\\
&5+2+1+1+1
\end{array}
\right\}\, &(4,1,0) \, &2
\end{align*}
\end{minipage}
%
\begin{minipage}{0.50\textwidth}
\small
\begin{align*}
& \mbox{ Partitions in $Q$ } &(l_{1}, l_{2},l_{3}) \, &\sharp\\
&\left\{
\begin{array}{lll}
&5+3+2\\
&6+3+1\\
&7+2+1
\end{array}
\right\}\, &(1,1,1) \, &3\\
&\left\{
\begin{array}{ll}
&5+2+2+1\\
&6+2+1+1
\end{array}
\right\}\, &(2,2,0) \, &2\\
&\left\{
\begin{array}{ll}
&5+3+1+1\\
&7+1+1+1
\end{array}
\right\}\, &(3,0,1) \, &2\\
&\left\{
\begin{array}{ll}
&5+2+1+1+1\\
&6+1+1+1+1
\end{array}
\right\}\, &(4,1,0) \, &2
\end{align*}
\end{minipage}

\begin{minipage}{0.43\textwidth}
\small
\begin{align*}
&\left\{
\begin{array}{l}
4+4+2
\end{array}
\right\}\,\quad\quad\quad\quad\quad\quad\quad &(0,2,2)  &1\\
&\left\{
\begin{array}{l}
4+3+3
\end{array}
\right\}\, &(1,0,3) \, &1\\
&\left\{
\begin{array}{l}
5+4+1
\end{array}
\right\}\, &(1,3,1) \, &1\\
&\left\{
\begin{array}{l}
5+3+2
\end{array}
\right\}\, &(2,1,2) \, &1\\
&\left\{
\begin{array}{l}
6+4
\end{array}
\right\}\, &(2,4,0) \, &1\\
&\left\{
\begin{array}{l}
6+3+1
\end{array}
\right\}\, &(3,2,1) \, &1\\
&\left\{
\begin{array}{l}
6+2+2
\end{array}
\right\}\, &(4,0,2) \, &1\\
&\left\{
\begin{array}{l}
7+3
\end{array}
\right\}\, &(4,3,0) \, &1\\
&\left\{
\begin{array}{l}
7+2+1
\end{array}
\right\}\, &(5,1,1) \, &1\\
&\left\{
\begin{array}{l}
8+2
\end{array}
\right\}\, &(6,2,0) \,&1\\
&\left\{
\begin{array}{l}
8+1+1
\end{array}
\right\}\, &(7,0,1) \, &1\\
&\left\{
\begin{array}{l}
9+1
\end{array}
\right\}\, &(8,1,0) \, &1
\end{align*}
\end{minipage}
%
\begin{minipage}{0.50\textwidth}
\small
\begin{align*}
\vspace{0.5cm}
&\left\{
\begin{array}{l}
3+3+2+2
\end{array}
\right\}\, &(0,2,2) \, &1\\
&\left\{
\begin{array}{l}
3+3+3+1
\end{array}
\right\}\, &(1,0,3) \, &1\\
&\left\{
\begin{array}{l}
3+2+2+2+1
\end{array}
\right\}\, &(1,3,1) \, &1\\
&\left\{
\begin{array}{l}
3+3+2+1+1
\end{array}
\right\}\, &(2,1,2) \, &1\\
&\left\{
\begin{array}{l}
2+2+2+2+1+1
\end{array}
\right\}\, &(2,4,0) \, &1\\
&\left\{
\begin{array}{l}
3+2+2+1+1+1
\end{array}
\right\}\, &(3,2,1) \, &1\\
&\left\{
\begin{array}{l}
3+3+1+1+1+1
\end{array}
\right\}\, &(4,0,2) \, &1\\
&\left\{
\begin{array}{l}
2+2+2+1+1+1+1
\end{array}
\right\}\, &(4,3,0) \, &1\\
&\left\{
\begin{array}{l}
3+2+1+1+1+1+1
\end{array}
\right\}\, &(5,1,1) \, &1\\
&\left\{
\begin{array}{l}
2+2+1+1+1+1+1+1
\end{array}
\right\}\, &(6,2,0) \, &1\\
&\left\{
\begin{array}{l}
3+1+1+1+1+1+1+1
\end{array}
\right\}\, &(7,0,1) \, &1\\
&\left\{
\begin{array}{l}
2+1+1+1+1+1+1+1+1
\end{array}
\right\}\, &(8,1,0) \,&1
\end{align*}  
\end{minipage}



%%%%%%%%%%%%%%%%%%%%%%%%%%%
%%%%%%%%%%%%%%%%%%%%%%%%%%%%








%\end{example}

\section{ Main results and proofs}
In this section, we prove our conjecture is true for two family partitions classes, which is the main results of this paper.
\begin{theorem}
The conjecture is true if $m=3$ and the type of partitions is $(\Sigma, 2)$ or $(2, \Sigma)$. $\Sigma$ can be any positive 
integer.
\end{theorem}
\begin{theorem}
For any $m\geq 3$, the conjecture is true if the type $(\Sigma_{1}, \Sigma_{2}, \dots, \Sigma_{m-1})$ satisfies the conditions:

$\bullet$ one $\Sigma$ is $1$,

$\bullet$ another $\Sigma$ can be any positive integer,

$\bullet$ other $m-3$ $\Sigma$ are zero.
Hence if $m=4$, $\Sigma>0$, this means that the conjecture is true for the following types:
$$
(\Sigma, 1, 0),\quad(\Sigma, 0, 1),\quad(0, \Sigma, 1),\quad(1, \Sigma, 0),\quad(1, 0,  \Sigma),\quad(0, 1,  \Sigma).
$$
\end{theorem}
In fact, for small $m$ and $\Sigma_{i}\leq 2$, we proved this conjecture is true, since the method for proving them
is the same as that is used to prove Theorem 3.1 and Theorem 3.2 but with less complication, so we omit these parts.
We should note that for general $n$,  most of partitions belong to $P$ have the types appearing in Theorem 3.1 and Theorem 3.2. For example, for $n=16$, the number of the partitions belong to $P$ is $89$, 19 of them have pure types, 56 of them have the types appearing in Theorem 3.1 and Theorem 3.2, only $14$ of them have other mixed types.

Our method for proving Theorem 3.1 and Theorem 3.2 is the comparison of the generating functions for partitions appearing in both sides of the conjecture and verifying both of them satisfying the same recurrences.
\begin{lemma}
Let $m\geq 2$, $1\leq i\leq m-1$. Let $Q_{i}$ be the set of partitions with all parts are congruent to $i$ modulo $m$, Let $f_{Q_{i}}(z,q)$ be the generating function for $Q_{i}$, $b_{i}(l; n)$ be the number of partitions of $n$ with $l$ parts and all parts $\equiv i\pmod m$.  Let $f_{Q}(z_{1}, z_{2}, \dots, z_{m-1},q)$ be the generating function of $Q$, and $b(l_{1}, l_{2}, \dots, l_{m-1}; n)$ be the number of partitions of $n$ with length type $(l_{1}, l_{2}, \dots, l_{m-1})$. That is
$$
f_{Q_{i}}(z,q)=\sum_{\lambda \in Q_{i}}z^{l(\lambda)} q^{|\lambda |}   =\sum_{l\geq 0,n\geq0}b_{i}(l;n)z^{l}q^{n}, \quad  l(\lambda)=\mbox{ length of } \lambda,
$$
and
$$
f_{Q}(z_{1}, z_{2}, \dots, z_{m-1},q)=\sum_{\lambda \in Q}z_{1}^{l_{1}(\lambda)}z_{2}^{l_{2}(\lambda)}\dots z_{m-1}^{l_{m-1}(\lambda)}q^{|\lambda |}=\sum_{l_{i}, n\geq 0} b(l_{1}, l_{2}, \dots, l_{m-1}; n) z_{1}^{l_{1}}z_{2}^{l_{2}}\dots z_{m-1}^{l_{m-1}}q^{n}.
$$
Then the generating function of $Q$ satisfies
$$
f_{Q}(z_{1}, z_{2}, \dots, z_{m-1},q)=\prod_{i=1}^{m-1}f_{Q_{i}}(z_{i}, q).
$$
Equivalently, 
$$
b(l_{1}, l_{2}, \dots, l_{m-1}; n)=\sum_{n_{i}\geq 0, n_{1}+\dots n_{m-1}=n} b_{1}(l_{1};n_{1})b_{2}(l_{2};n_{2})\dots b_{m-1}(l_{m-1}; n_{m-1}).
$$
\end{lemma}
\noindent{\textit{Proof.}}
This is an application of basic counting principle. In order to construct a partition of $n$ with length type $(l_{1}, l_{2}, \dots, l_{m-1})$, we firstly decomposition $n$ into $m-1$ positive integers $n_{1}, n_{2},\dots, n_{m-1}$, for  each $n_{i}$, choose 
a partition of $n_{i}$ with all parts $\equiv i \pmod m$ and $l_{i}$ parts. Put all parts of these m-1 partitions together, we get a partition of $n$ with length type $(l_{1}, l_{2}, \dots, l_{m-1})$. When considering all partitions of $n_{i}$ with the given conditions and  all possible decompositions of $n$, we get all partitions of $n$ with given the length type $(l_{1}, l_{2}, \dots, l_{m-1})$.

\noindent From this lemma and the facts that the function
$$ \frac{q^{2i}}{(1-q^{3})(1-q^{6})} \quad (i=1, 2) $$ 
generates all partitions with exactly two parts and each part $\equiv i\pmod 3$ and 
$\frac{q^{j}}{1-q^{m}}(1\leq j\leq m-1)$ generates all partitions consisting only one part and this part is congruent to $j$ modulo $m$,
we get the following results:
\begin{align}
&\sum_{l, n\geq 0}b(l, 2; n)z^{l}q^{n}=\left(\sum_{l, n\geq 0}b(l,0; n)z^{l}q^{n} \right)\frac{q^{4}}{(1-q^{3})(1-q^{6})},\label{3.1}\\
&\sum_{l, n\geq 0}b(2, l; n)z^{l}q^{n}=\left(\sum_{l, n\geq 0}b(0, l; n)z^{l}q^{n} \right)\frac{q^{2}}{(1-q^{3})(1-q^{6})},\label{3.2}\\
&\sum_{l, n\geq 0}\tilde{b}_{i; j}(l;n)z^{l}q^{n}=\left(\sum_{l, n\geq 0}b(0,\dots,0,l,0,\dots, 0;n)z^{l}q^{n} \right)\frac{q^{j}}{1-q^{m}}.\label{3.3}
\end{align}
Where $1\leq j\leq m-1, j\not\neq i$ and $\tilde{b}_{i;j}(l;n)=b(0,0,\dots, 0, l ,0,\dots,0,1,0,\dots,0; n)$, which is the number of partitions of $n$ with exactly $l$ parts
$\equiv i \pmod m$ and one part $\equiv j \pmod m$.

Now we prove Theorem 3.1 and Theorem 3.2.

\noindent {\it Proof of Theorem 3.1.}

We prove Theorem 3.1 for the type $(\Sigma, 2)$, and the proof for the case of the type $(2, \Sigma)$ is similar. Let $a(\Sigma_{1}, \Sigma_{2};  n)$ denote the number of partitions of n with the alternating sum type
$(\Sigma_{1},  \Sigma_{2})$. Theorem 3.1 is equivalent to the claim:  for any $\Sigma \geq 0$, $a(\Sigma, 2;  n)=b(\Sigma, 2; n)$. Since $a(\Sigma, 0;  n)=b(\Sigma, 0; n)$, the conjecture for the pure type $(\Sigma,0)$ , now is Theorem 1.1. So it is only to prove the generating function of $a(\Sigma, 2; n)$ satisfies the following identity by comparing with the identity (\ref{3.1}):
\begin{align}
\sum_{\Sigma, n\geq 0}a(\Sigma, 2;n)z^{\Sigma}q^{n}=\left(\sum_{\Sigma, n\geq 0}a(\Sigma, 0 ;n)z^{\Sigma}q^{n} \right)\frac{q^{4}}{(1-q^{3})(1-q^{6})}.
\end{align}
Let $\lambda$ be a partition with parts repeated times $\leq 2$ and the alternating sum type $(\Sigma, 2)$. The basic units in 
$\lambda$ have two cases:

$\bullet$ \textbf{ Case A.}\label{Case A} The basic units  consist in forms:
$$
\lambda_{1} >\lambda_{2}=\lambda_{3} \quad \mbox{ and  two forms}\quad \lambda_{1}^{\prime}\geq \lambda_{2}^{\prime}\overset{1}{>}\lambda_{3}^{\prime}.
$$ 
Where $\lambda_{2}^{\prime}\overset{1}{>}\lambda_{3}^{\prime}$ means $\lambda_{2}^{\prime}-\lambda_{3}^{\prime}=1$.  Here is an example with four basic units:
$$
10>8=8 >{\color{blue}7=7>6}     =6>5=5 >{\color{blue}2>2>1},$$ the first  unit and third unit are the form $\lambda_{1} >\lambda_{2}=\lambda_{3} $, the second unit and fourth unit (with blue colour) are the form $\lambda_{1}^{\prime}\geq \lambda_{2}^{\prime}\overset{1}{>}\lambda_{3}^{\prime}.$

$\bullet$ \textbf{ Case B.} \label{Case B}The basic units consist in forms:
$$
\lambda_{1} >\lambda_{2}=\lambda_{3} \quad \mbox{ and one form}\quad \lambda_{1}^{\prime}\geq \lambda_{2}^{\prime}\overset{2}{>}\lambda_{3}^{\prime}.
$$
Where $\lambda_{2}^{\prime}\overset{2}{>}\lambda_{3}^{\prime}$ means $\lambda_{2}^{\prime}-\lambda_{3}^{\prime}=2$. We give an example for this case,
$
8>7=7>6>4=4>{\color{blue}3>2},
$
where the first unit is $\lambda_{1}=8, \lambda_{2}=\lambda_{3}=7$,  the second unit is $\lambda_{4}=6, \lambda_{5}=\lambda_{6}=4$, and the third unit is ${\color{blue}\lambda_{7}=3,\lambda_{8}=2, \lambda_{9}=0}$.



We call the types $\lambda_{1}^{\prime}\geq \lambda_{2}^{\prime}\overset{2}{>}\lambda_{3}^{\prime}$ and $\lambda_{1}^{\prime}\geq \lambda_{2}^{\prime}\overset{1}{>}\lambda_{3}^{\prime}$ to be special units.

We firstly consider \textbf{Case A}.
\begin{lemma}\label{Lemma 3.4}%%%%%%%%%%%%%%%%%%%%%%%%%  lemma 3.4
 The generating function for the number of partitions of fixed length $3n+3$ with alternating sum type $(\Sigma, 2)$ and belonging to the {\textbf{Case A}} is the sum of the following $n$ terms, where $n\geq 1$.




\begin{align}\tag{0}\label{3.4-0}
&\frac{z^{n+1}q^{3n^{2}+7n+4}}{(zq;q^{3})_{n+1}(q^{3};q^{3})_{n+1}}\nonumber
\left( \sum_{k=1} ^{n-1}\left(\frac{1}{z^{2}q^{3k+1}}+\frac{1-q^{3k}}{zq^{3k}}\right)(n\geq2)+\left(\frac{q^{2}}{z^{2}}+\frac{q^{3}(1-q^{3n})}{z}\right)(n\geq 1)  \right)\nonumber,
\end{align}



\begin{equation}\tag{1}\label{3.4-1}
\frac{z^{n+1}q^{3n^{2}+7n+4}}{(zq;q^{3})_{n+1}(q^{3};q^{3})_{n+1}}\left( \sum_{k=1} ^{n-2}\frac{1}{z^{2}q^{6k+4}}(n\geq 3)+\frac{1}{z^{2}q^{3n-5}}(n\geq 2)\right),
\end{equation}



\begin{equation}\tag{2}
\frac{z^{n+1}q^{3n^{2}+7n+4}}{(zq;q^{3})_{n+1}(q^{3};q^{3})_{n+1}}\left( \sum_{k=1} ^{n-3}\frac{1}{z^{2}q^{6k+7}}(n\geq 4)+\frac{1}{z^{2}q^{3n-8}}(n\geq 3)\right),
\end{equation}

\begin{equation}\nonumber
\dots
\end{equation}


\begin{equation}\tag{n-2}
\frac{z^{n+1}q^{3n^{2}+7n+4}}{(zq;q^{3})_{n+1}(q^{3};q^{3})_{n+1}}\left( \sum_{k=1} ^{n-(n-1)}\frac{1}{z^{2}q^{6k+3n-5}}+\frac{1}{z^{2}q^{4}}\right),
\end{equation}

\begin{equation}\tag{n-1}\label{3.4-n-1}
\frac{z^{n+1}q^{3n^{2}+7n+4}}{(zq;q^{3})_{n+1}(q^{3};q^{3})_{n+1}}\frac{1}{z^{2}q^{}}.
\end{equation}

\end{lemma}
 Here we write them as a list  so that  we can easily see how each term of them corresponds to the following analysis and the inequalities on $n$ in brackets means the corresponding terms requiring $n$ satisfying  this condition. Moreover we use the standard $q$-Pochhammer symbol notations  for simplicity defined by 
 \begin{equation*}
(a;q)_n := \prod_{i=0}^{n-1} (1-aq^{i}) \text{  and  } (a;q)_\infty :=\lim_{L\rightarrow\infty} (a;q)_n.
\end{equation*}
But in the proof of our main results, we sometimes choose product forms instead of $q$-Pochhammer symbol so that readers can easily track our transformations.


\begin{lemma}\label{Lemma 3.5}%%%%%%%%%%%%%%%% lemma 3.5
The generating function  for the number of partitions of fixed length $3n+2$ with the alternating sum type $(\Sigma, 2)$ and belonging to the \textbf{Case A} is the sum of the following $n$ terms.
\begin{equation}\tag{0}\label{3.5-0}
\frac{z^{n+1}q^{3n^{2}+4n+1}}{(zq;q^{3})_{n+1}(q^{3};q^{3})_{n}} \left(\frac{q^{2}}{z^{2}}+\frac{q^{3}(1-q^{3n})}{z}\right)(n\geq 1),
\end{equation}



\begin{equation}\tag{1}\label{3.5-1}
\frac{z^{n+1}q^{3n^{2}+4n+1}}{(zq;q^{3})_{n+1}(q^{3};q^{3})_{n}} \frac{1}{z^{2}q^{3n-5}}(n\geq 2),
\end{equation}



\begin{equation}\tag{2}
\frac{z^{n+1}q^{3n^{2}+4n+1}}{(zq;q^{3})_{n+1}(q^{3};q^{3})_{n}}\frac{1}{z^{2}q^{3n-8}}(n\geq 3),
\end{equation}

\begin{equation*}
\dots
\end{equation*}


\begin{equation}\tag{n-2}
\frac{z^{n+1}q^{3n^{2}+4n+1}}{(zq;q^{3})_{n+1}(q^{3};q^{3})_{n}}\frac{1}{z^{2}q^{4}},
\end{equation}

\begin{equation}\tag{n-1}\label{3.5-n-1}
\frac{z^{n+1}q^{3n^{2}+4n+1}}{(zq;q^{3})_{n+1}(q^{3};q^{3})_{n}}\frac{1}{z^{2}q^{}}.
\end{equation}


\end{lemma}
%%%%%%%%%%%%%%%%
%%%%%%%%%%%%%%%%%      Lemma3.6
\begin{lemma}\label{Lemma 3.6}
The generating function for the number of partitions of fixed length $3n+1$ with alternating sum type $(\Sigma, 2)$ and belonging to the \textbf{Case A} is the sum of the following $n-1$ terms. 
\begin{align}\tag{0}\label{3.6-0}
\frac{z^{n+1}q^{3n^{2}+4n+1}}{(zq;q^{3})_{n+1}(q^{3};q^{3})_{n}}
\sum_{k=1} ^{n-1}\left(\frac{1}{z^{2}q^{3k+1}}+\frac{1-q^{3k}}{zq^{3k}}\right)(n\geq2),
\end{align}



\begin{equation}\tag{1}
\frac{z^{n+1}q^{3n^{2}+4n+1}}{(zq;q^{3})_{n+1}(q^{3};q^{3})_{n}}\sum_{k=1} ^{n-2}\frac{1}{z^{2}q^{6k+4}}(n\geq 3),
\end{equation}



\begin{equation}\tag{2}
\frac{z^{n+1}q^{3n^{2}+4n+1}}{(zq;q^{3})_{n+1}(q^{3};q^{3})_{n}}\sum_{k=1} ^{n-3}\frac{1}{z^{2}q^{6k+7}}(n\geq 4),
\end{equation}

\begin{equation*}
\dots
\end{equation*}


\begin{equation}\tag{n-2}\label{3.6-n-2}
\frac{z^{n+1}q^{3n^{2}+4n+1}}{(zq;q^{3})_{n+1}(q^{3};q^{3})_{n}}\sum_{k=1} ^{n-(n-1)}\frac{1}{z^{2}q^{6k+3n-5}}.
\end{equation}



\end{lemma}
\noindent {\it Proof of Lemma 3.4.}
Let $\lambda$ be a partition of length $3n+3$ and its basic units belonging to the \textbf{Case A.} Note that $\lambda$ has $n+1$ basic units, two of them are special units. We classify such partitions into $n$ classes by the distance of two special units. Here the distance $0$ means that there is no basic unit of the type $\lambda_{1}>\lambda_{2}=\lambda_{3}$ between the two special units. The distance $1$ means there is only one basic unit of type $\lambda_{1}>\lambda_{2}=\lambda_{3}$ between the two special units, etc. For example, the partition
$$
9>8=8>{\color{blue}7=7>6}>5>4=4>{\color{blue}3>2>1},
$$
has the distance between special units (the color is blue) $1$. By a partition of class $d$, $0\leq d\leq n-1$, it means the distance between special units is $d$. We consider each class as follows.

$\bullet$ {\bf Distance $0.$} Let $\lambda$ be such a partition, if the two special units are not the last two basic units,
 then it has the form:
\begin{align*}
&\lambda_{1}>\lambda_{2}=\lambda_{3}>\dots>\lambda_{3k-5}>\lambda_{3k-4}=\lambda_{3k-3}\\
>&{\color{blue}\lambda_{3k-2}\geq\lambda_{3k-1}\overset{1}{>}\lambda_{3k}>\lambda_{3k+1}\geq \lambda_{3k+2}\overset{1}{>}\lambda_{3k+3}}\\
\geq &\lambda_{3k+4}>\lambda_{3k+5}=\lambda_{3k+6}>\dots >\lambda_{3n+1}>\lambda_{3n+2}=\lambda_{3n+3},
\end{align*}
 or the form:
\begin{align*}
&\lambda_{1}>\lambda_{2}=\lambda_{3}>\dots>\lambda_{3k-5}>\lambda_{3k-4}=\lambda_{3k-3}\\
>&{\color{blue}\lambda_{3k-2}\geq\lambda_{3k-1}\overset{1}{>}\lambda_{3k}=\lambda_{3k+1}> \lambda_{3k+2}\overset{1}{>}\lambda_{3k+3}}\\
\geq &\lambda_{3k+4}>\lambda_{3k+5}=\lambda_{3k+6}>\dots >\lambda_{3n+1}>\lambda_{3n+2}=\lambda_{3n+3}.
\end{align*}
Where $\lambda_{3n+3}>0$ and $1\leq k\leq n-1$. By considering conjugates of such partitions, the standard partition analysis gives the generating function for the first case is
\begin{align}
 &\frac{zq}{1-zq}\frac{q^{3}}{1-q^{3}}\dots\frac{zq^{3k-5}}{1-zq^{3k-5}}\frac{q^{3k-3}}{1-q^{3k-3}}
 \frac{1}{1-zq^{3k-2}}\cdot q^{3k-1}\cdot\frac{q^{3k}}{1-q^{3k}}\nonumber\\
 &\cdot\frac{1}{1-zq^{3k+1}}\cdot q^{3k+2}\frac{1}{1-q^{3k+3}}
\cdot\frac{zq^{3k+4}}{1-zq^{3k+4}}\dots\frac{zq^{3n+1}}{1-zq^{3n+1}}\frac{q^{3n+3}}{1-q^{3n+3}}\nonumber\\
= &\frac{zq}{1-zq}\frac{q^{3}}{1-q^{3}} \frac{zq^{4}}{1-zq^{4}}\dots\frac{zq^{3n+1}}{1-zq^{3n+1}} \frac{q^{3n+3}}{1-q^{3n+3}}\cdot\frac{q^{2}}{z^{2}}\cdot\frac{1}{q^{3k+3}}\nonumber\\
=&\frac{z^{n+1}q^{3n^{2}+7n+4}}{(zq;q^{3})_{n+1}(q^{3};q^{3})_{n+1}}\frac{1}{z^{2}q^{3k+1}},\label{3.5}
\end{align}
and the generating function for the second case is
\begin{align}
 &\frac{zq}{1-zq}\frac{q^{3}}{1-q^{3}}\dots\frac{zq^{3k-5}}{1-zq^{3k-5}}\frac{q^{3k-3}}{1-q^{3k-3}}
 \frac{1}{1-zq^{3k-2}}\cdot q^{3k-1}\nonumber\\
& \cdot\frac{zq^{3k+1}}{1-zq^{3k+1}}\cdot q^{3k+2}\frac{1}{1-q^{3k+3}}
\cdot\frac{zq^{3k+4}}{1-zq^{3k+4}}\dots\frac{zq^{3n+1}}{1-zq^{3n+1}}\frac{q^{3n+3}}{1-q^{3n+3}}\nonumber\\
=& \frac{zq}{1-zq}\frac{q^{3}}{1-q^{3}} \frac{zq^{4}}{1-zq^{4}}\dots\frac{zq^{3n+1}}{1-zq^{3n+1}} \frac{q^{3n+3}}{1-q^{3n+3}}\cdot\frac{q^{3}(1-q^{3k})}{z}\cdot\frac{1}{q^{3k+3}}\nonumber\\
=&\frac{z^{n+1}q^{3n^{2}+7n+4}}{(zq;q^{3})_{n+1}(q^{3};q^{3})_{n+1}}\frac{1-q^{3k}}{zq^{3k}}.\label{3.6}
 \end{align}

If the special units are the last two basic units, then it has the form:
\begin{align*}
&\lambda_{1}>\lambda_{2}=\lambda_{3}>\dots>\lambda_{3n-5}>\lambda_{3n-4}=\lambda_{3n-3}\\
>&{\color{blue}\lambda_{3n-2}\geq\lambda_{3n-1}\overset{1}{>}\lambda_{3n}>\lambda_{3n+1}\geq \lambda_{3n+2}\overset{1}{>}\lambda_{3n+3}},\quad
 (\lambda_{3n+3}>0),
\end{align*} or the form:
\begin{align*}
&\lambda_{1}>\lambda_{2}=\lambda_{3}>\dots>\lambda_{3n-5}>\lambda_{3n-4}=\lambda_{3n-3}\\
>&{\color{blue}\lambda_{3n-2}\geq\lambda_{3n-1}\overset{1}{>}\lambda_{3n}=\lambda_{3n+1}> \lambda_{3n+2}\overset{1}{>}\lambda_{3n+3}},\quad
(\lambda_{3n+3}>0).
\end{align*}
By considering conjugates of such partitions, the standard partition analysis gives the generating function for the first case is
\begin{align}
 \frac{zq}{1-zq}\frac{q^{3}}{1-q^{3}}\dots\frac{zq^{3n-5}}{1-zq^{3n-5}}\frac{q^{3n-3}}{1-q^{3n-3}}
 \frac{1}{1-zq^{3n-2}}\cdot &q^{3n-1}\cdot\frac{q^{3n}}{1-q^{3n}}\cdot\frac{1}{1-zq^{3n+1}}\cdot q^{3n+2}\frac{q^{3n+3}}{1-q^{3n+3}}\nonumber\\
= \frac{zq}{1-zq}\frac{q^{3}}{1-q^{3}} \frac{zq^{4}}{1-zq^{4}}\dots&\frac{zq^{3n+1}}{1-zq^{3n+1}} \frac{q^{3n+3}}{1-q^{3n+3}}\cdot\frac{q^{2}}{z^{2}}\nonumber\\
=\frac{z^{n+1}q^{3n^{2}+7n+4}}{(zq;q^{3})_{n+1}(q^{3};q^{3})_{n+1}}&\frac{q^{2}}{z^{2}}.\label{3.7}
\end{align}
And the generating function for the second case is
\begin{align}
 \frac{zq}{1-zq}\frac{q^{3}}{1-q^{3}}\dots\frac{zq^{3n-5}}{1-zq^{3n-5}}\frac{q^{3n-3}}{1-q^{3n-3}}
 &\frac{1}{1-zq^{3n-2}}\cdot q^{3n-1}\cdot\frac{zq^{3n+1}}{1-zq^{3n+1}}\cdot q^{3n+2}\frac{q^{3n+3}}{1-q^{3n+3}}\nonumber\\
= \frac{zq}{1-zq}\frac{q^{3}}{1-q^{3}} \frac{zq^{4}}{1-zq^{4}}\dots&\frac{zq^{3n+1}}{1-zq^{3n+1}} \frac{q^{3n+3}}{1-q^{3n+3}}\cdot\frac{q^{3}(1-q^{3n})}{z}\nonumber\\
=\frac{z^{n+1}q^{3n^{2}+7n+4}}{(zq;q^{3})_{n+1}(q^{3};q^{3})_{n+1}}&\frac{q^{3}(1-q^{3n})}{z}.\label{3.8}
 \end{align}
 
Add up (\ref{3.5}) and (\ref{3.6}), summing $k$ from $1$ to $n-1$ and then add them to (\ref{3.7}) and (\ref{3.8}), the sum is
\begin{align}\nonumber
&\frac{z^{n+1}q^{3n^{2}+7n+4}}{(zq;q^{3})_{n+1}(q^{3};q^{3})_{n+1}}\left( \sum_{k=1} ^{n-1}\left(\frac{1}{z^{2}q^{3k+1}}+\frac{1-q^{3k}}{zq^{3k}}\right)+\left(\frac{q^{2}}{z^{2}}+\frac{q^{3}(1-q^{3n})}{z}\right)  \right)\nonumber,
\end{align}
which is the term (\ref{3.4-0}) in Lemma 3.4. When special units are not the last basic unit, such a partition has at least three basic units, hence in the term
$$
\sum_{k=1} ^{n-1}\left(\frac{1}{z^{2}q^{3k+1}}+\frac{1-q^{3k}}{zq^{3k}}\right),
$$
$n$ should satisfy $n\geq 2$. The term $$\frac{q^{2}}{z^{2}}+\frac{q^{3}(1-q^{3n})}{z}$$ corresponds to the special units being
the last two basic units, it implies $\lambda$ has at least two basic units, hence $n\geq 1$. In the following cases, we have similar analysis on inequalities involving $n$. 

$\bullet$ {\bf Distance $1.$}
In this case, the two special  units must appear in the following form:
$$
{\color{blue}\lambda_{3k-2}  \geq \lambda_{3k-1} \overset{1}{>} \lambda_{3k}} \geq  \lambda_{3k+1} > \lambda_{3k+2} = \lambda_{3k+3} > {\color{blue} \lambda_{3k+4} \geq \lambda_{3k+5}\overset{1}{>}  \lambda_{3k+6}} ,
$$
if any special unit  does not  appear as the last basic unit, hence $k$ can take from $1$ to $n-2$, or the form
$$
{\color{blue}\lambda_{3n-5}  \geq \lambda_{3n-4} \overset{1}{>} \lambda_{3n-3}} \geq  \lambda_{3n-2} > \lambda_{3n-1} = \lambda_{3n} > {\color{blue} \lambda_{3n+1} \geq \lambda_{3n+2}\overset{1}{>}  \lambda_{3n+3}}.
$$
if one special unit is exactly the last basic unit. The generating function for the former is
\begin{align}
 &\frac{zq}{1-zq}\frac{q^{3}}{1-q^{3}}\dots\frac{zq^{3k-5}}{1-zq^{3k-5}}\frac{q^{3k-3}}{1-q^{3k-3}}
 \frac{1}{1-zq^{3k-2}}\cdot q^{3k-1}\cdot\frac{1}{1-q^{3k}}\cdot\nonumber\\
 &\frac{zq^{3k+1}}{1-zq^{3k+1}}\cdot \frac{q^{3k+3}}{1-q^{3k+3}}
 \frac{1}{1-zq^{3k+4}}\cdot q^{3k+5}\cdot\frac{1}{1-q^{3k+6}}
\cdot\frac{zq^{3k+7}}{1-zq^{3k+7}}\dots\frac{zq^{3n+1}}{1-zq^{3n+1}}\frac{q^{3n+3}}{1-q^{3n+3}}\nonumber\\
&= \frac{zq}{1-zq}\frac{q^{3}}{1-q^{3}} \frac{zq^{4}}{1-zq^{4}}\dots\frac{zq^{3n+1}}{1-zq^{3n+1}} \frac{q^{3n+3}}{1-q^{3n+3}}\cdot\frac{1}{z^{2}q^{6k+4}} \quad (n \geq 3)\nonumber\\
&=\frac{z^{n+1}q^{3n^{2}+7n+4}}{(zq;q^{3})_{n+1}(q^{3};q^{3})_{n+1}}\frac{1}{z^{2}q^{6k+4}}, \quad(n\geq3),\label{3.9}
\end{align}
 and for the later is
\begin{align}
 &\frac{zq}{1-zq}\frac{q^{3}}{1-q^{3}}\dots\frac{zq^{3n-7}}{1-zq^{3n-7}}\frac{q^{3n-6}}{1-q^{3n-6}}
 \frac{1}{1-zq^{3n-5}}\cdot q^{3n-4}\cdot\frac{1}{1-q^{3n-3}}\cdot\nonumber\\
 &\frac{zq^{3n-2}}{1-zq^{3n-2}}\cdot \frac{q^{3n}}{1-q^{3n}}\cdot\frac{1}{1-zq^{3n+1}}\cdot q^{3n+2}\cdot\frac{q^{3n+3}}{1-q^{3n+3}}\nonumber\\
&= \frac{zq}{1-zq}\frac{q^{3}}{1-q^{3}} \frac{zq^{4}}{1-zq^{4}}\dots\frac{zq^{3n+1}}{1-zq^{3n+1}} \frac{q^{3n+3}}{1-q^{3n+3}}\cdot\frac{1}{z^{2}q^{3n-5}}  \quad( n\geq 2)\nonumber\\
&=\frac{z^{n+1}q^{3n^{2}+7n+4}}{(zq;q^{3})_{n+1}(q^{3};q^{3})_{n+1}}\frac{1}{z^{2}q^{3n-5}}, \quad(n\geq2).\label{3.10}
\end{align}
Sum (\ref{3.9}) for $k$ from $1$ to $n-2$ and then add it to (\ref{3.10}), we get the term
\begin{equation*}
\frac{z^{n+1}q^{3n^{2}+7n+4}}{(zq;q^{3})_{n+1}(q^{3};q^{3})_{n+1}}\left( \sum_{k=1} ^{n-2}\frac{1}{z^{2}q^{6k+4}}(n\geq 3)+\frac{1}{z^{2}q^{3n-5}}(n\geq 2)\right),
\end{equation*}
 which is the term (\ref{3.4-1}) in Lemma 3.4.
 
 
$\bullet$ {\bf Distance $d$ satisfies  $2\leq d\leq n-2$.}
 The two special units must appear in the form:
$$
{\color{blue}\lambda_{3k-2}  \geq \lambda_{3k-1} \overset{1}{>} \lambda_{3k}} \geq  \lambda_{3k+1} > \lambda_{3k+2} = \lambda_{3k+3} > \dots=\lambda_{3k+3d}> {\color{blue}\lambda_{3k+3d+1}\geq\lambda_{3k+3d+2} \overset{1}{>} \lambda_{3k+3d+3}},
$$
corresponding to the last special unit  being not the last basic unit, hence $k$ can take from $1$ to $n-d-1$ for fixed $d$, or in the form
$$
{\color{blue}\lambda_{3n-3d-2}  \geq \lambda_{3n-3d-1} \overset{1}{>} \lambda_{3n-3d}} \geq  \lambda_{3n-3d+1} >\dots > {\color{blue} \lambda_{3n+1} \geq \lambda_{3n+2}\overset{1}{>}  \lambda_{3n+3}},
$$
corresponding to one special unit exactly being the last basic unit. The generating function for the former is
\begin{align}
 &\frac{zq}{1-zq}\frac{q^{3}}{1-q^{3}}\dots\frac{zq^{3k-5}}{1-zq^{3k-5}}\frac{q^{3k-3}}{1-q^{3k-3}}
 \frac{1}{1-zq^{3k-2}}\cdot q^{3k-1}\cdot\frac{1}{1-q^{3k}}\cdot
 \frac{zq^{3k+1}}{1-zq^{3k+1}}\dots\frac{zq^{3k+3d-2}}{1-q^{3k+3d-2}}\nonumber\\
 &\frac{q^{3k+3d}}{1-q^{3k+3d}}
 \frac{1}{1-zq^{3k+3d+1}}\cdot q^{3k+3d+2}\cdot\frac{1}{1-q^{3k+3d+3}}
\cdot\frac{zq^{3k+3d+4}}{1-zq^{3k+3d+4}}\dots\frac{zq^{3n+1}}{1-zq^{3n+1}}\frac{q^{3n+3}}{1-q^{3n+3}}\nonumber\\
&= \frac{zq}{1-zq}\frac{q^{3}}{1-q^{3}} \frac{zq^{4}}{1-zq^{4}}\dots\frac{zq^{3n+1}}{1-zq^{3n+1}} \frac{q^{3n+3}}{1-q^{3n+3}}\cdot\frac{1}{z^{2}q^{6k+3d+1}}  \quad (2 \leq d \leq n-2, 1\leq k \leq n-d-1)\nonumber\\
&=\frac{z^{n+1}q^{3n^{2}+7n+4}}{(zq;q^{3})_{n+1}(q^{3};q^{3})_{n+1}}\frac{1}{z^{2}q^{6k+3d+1}},              \quad(2\leq d\leq n-2,1\leq k\leq n-d-1).\label{3.11}
\end{align}
 And for the later is
\begin{align}
 &\frac{zq}{1-zq}\frac{q^{3}}{1-q^{3}}\dots\frac{q^{3n-3d-3}}{1-q^{3n-3d-3}}
 \frac{1}{1-zq^{3n-3d-2}}\cdot q^{3n-3d-1}\cdot\frac{1}{1-q^{3n-3d}}\cdot\nonumber\\
 &\frac{zq^{3n-3d+1}}{1-zq^{3n-3d+1}}\dots \frac{q^{3n}}{1-q^{3n}}\cdot\frac{1}{1-zq^{3n+1}}\cdot q^{3n+2}\cdot\frac{q^{3n+3}}{1-q^{3n+3}}\nonumber\\
&= \frac{zq}{1-zq}\frac{q^{3}}{1-q^{3}} \frac{zq^{4}}{1-zq^{4}}\dots\frac{zq^{3n+1}}{1-zq^{3n+1}} \frac{q^{3n+3}}{1-q^{3n+3}}\cdot\frac{1}{z^{2}q^{3n-3d-2}}  \quad( 2 \leq d \leq n-1)\nonumber\\
&=\frac{z^{n+1}q^{3n^{2}+7n+4}}{(zq;q^{3})_{n+1}(q^{3};q^{3})_{n+1}}\frac{1}{z^{2}q^{3n-3d-2}}. \quad(2 \leq d\leq n-1)\label{3.12}
\end{align}
For each $2\leq d\leq n-2$, sum (\ref{3.11}) for $k$ from $1$ to $n-d-1$ and then add up the term in (\ref{3.12}) corresponding to $d$, we get the term (d) in Lemma 3.4.

$\bullet$ {\bf Distance $n-1.$} 
This is the last case of Lemma 3.4, it corresponds to the distance $n-2$ being maximal. Hence one special unit is the first unit and another special unit is the last basic unit. $\lambda$ has the form:
$$
{\color{blue}\lambda_{1}  \geq \lambda_{2} \overset{1}{>} \lambda_{3}} \geq  \lambda_{4} >\dots > {\color{blue} \lambda_{3n+1} \geq \lambda_{3n+2}\overset{1}{>}  \lambda_{3n+3}}.
$$
The generating function is
\begin{align*}
 &\frac{1}{1-zq}\cdot q^{2}\frac{1}{1-q^{3}}
 \frac{zq^{4}}{1-zq^{4}}>\dots >\frac{q^{3n}}{1-q^{3n}}\frac{1}{1-zq^{3n+1}}\cdot q^{3n+2}\cdot\frac{q^{3n+3}}{1-q^{3n+3}}\\\nonumber
&= \frac{zq}{1-zq}\frac{q^{3}}{1-q^{3}} \frac{zq^{4}}{1-zq^{4}}\dots\frac{zq^{3n+1}}{1-zq^{3n+1}} \frac{q^{3n+3}}{1-q^{3n+3}}\cdot\frac{1}{z^{2}q^{}}\\\nonumber
&=\frac{z^{n+1}q^{3n^{2}+7n+4}}{(zq;q^{3})_{n+1}(q^{3};q^{3})_{n+1}}\frac{1}{z^{2}q}.\nonumber
\end{align*}
Which is the term (\ref{3.4-n-1}) in Lemma 3.4. We complete the proof of  Lemma 3.4.



\noindent {\it Proof of Lemma 3.5.}

The analysis above holds for the case of $\lambda$ has length $3n+2$, except that in each case ($0\leq d \leq n-1$, $d=$ distance), one special unit must be the last unit and $\lambda_{3n+2}=1$ .

We use the case of distance $0$ to illustrate them. $\lambda$ has the form:
\begin{align*}
\lambda_{1}>\lambda_{2}=\lambda_{3}>\dots
>{\color{blue}\lambda_{3n-2}\geq\lambda_{3n-1}\overset{1}{>}\lambda_{3n}>\lambda_{3n+1}\geq \lambda_{3n+2}}, (\lambda_{3n+2}=1),
\end{align*} 
or the form:
\begin{align*}
\lambda_{1}>\lambda_{2}=\lambda_{3}>\dots
>{\color{blue}\lambda_{3n-2}\geq\lambda_{3n-1}\overset{1}{>}\lambda_{3n}=\lambda_{3n+1}> \lambda_{3n+2}}, (\lambda_{3n+2}=1).
\end{align*} 
The generating function for the former is 
\begin{align}
& \frac{zq}{1-zq}\frac{q^{3}}{1-q^{3}}\dots\frac{q^{3n-3}}{1-q^{3n-3}}
 \frac{1}{1-zq^{3n-2}} q^{3n-1}\frac{q^{3n}}{1-q^{3n}}\frac{1}{1-zq^{3n+1}} q^{3n+2}\nonumber\\
=& \frac{zq}{1-zq}\frac{q^{3}}{1-q^{3}} \frac{zq^{4}}{1-zq^{4}}\dots\frac{zq^{3n+1}}{1-zq^{3n+1}} \frac{q^{2}}{z^{2}}\nonumber\\
=&\frac{z^{n+1}q^{3n^{2}+4n+1}}{(zq;q^{3})_{n+1}(q^{3};q^{3})_{n}}\frac{q^{2}}{z^{2}}.\label{3.13}
\end{align}
The generating function for the later is
\begin{align}
& \frac{zq}{1-zq}\frac{q^{3}}{1-q^{3}}\dots\frac{q^{3n-3}}{1-q^{3n-3}}
 \frac{1}{1-zq^{3n-2}} q^{3n-1}\frac{zq^{3n+1}}{1-zq^{3n+1}}q^{3n+2}\nonumber\\
=& \frac{zq}{1-zq}\frac{q^{3}}{1-q^{3}} \frac{zq^{4}}{1-zq^{4}}\dots\frac{zq^{3n+1}}{1-zq^{3n+1}} \frac{q^{3}(1-q^{3n})}{z}\nonumber\\
=&\frac{z^{n+1}q^{3n^{2}+4n+1}}{(zq;q^{3})_{n+1}(q^{3};q^{3})_{n}}\frac{q^{3}(1-q^{3n})}{z}.\label{3.14}
 \end{align} 
Add up (\ref{3.13}) and (\ref{3.14}), we get the term (\ref{3.5-0}) in Lemma 3.5. Such a partition has at least two basic units, hence $n\geq 1.$  By the similar analysis in Lemma (3.4), we  find  that for the cases of distance from $1$ to $n-1$, partitions with length $3n+2$ and one special unit is the last unit can be obtained from  partitions with length $3n+3$  and one special unit is the last basic unit by omitting the part $\lambda_{3n+3}$ and letting $\lambda_{3n+2}=1$. Therefore, in each case, the generating function for partitions  with length $3n+2$ and one special unit being the last basic unit is the same as that of the partitions of length $3n+3$ with one special unit being the last basic unit except without the factor 
$$\frac{q^{3n+3}}{1-q^{3n+3}}.$$
They correspond to the terms 
(\ref{3.5-1}) to (\ref{3.5-n-1}) in Lemma 3.5.

We complete the proof of Lemma 3.5.




\noindent {\it Proof of Lemma 3.6.}

Since such a partition has length $3n+1$,  hence the last basic unit can not be a special unit. Therefore, the range of distance 
between two special units is from $1$ to $n-2$. As the analysis of the partitions with length $3n+3$, we find in each case
(distance from $0$ to $n-2$), partitions with length $3n+1$ and the last basic unit is not a special unit can be obtained from
the partitions with length $3n+3$ and the last basic unit being not a special unit by omitting the parts $\lambda_{3n+3}$ and $\lambda_{3n+2}$. Therefore, the generating function, for each $0\leq d\leq n-2$,  for partitions of length $3n+1$ and  the last basic unit being not a special unit can be obtained from the generating function for partitions with length $3n+3$ and the last basic unit being not a special unit by omitting the factor $$\frac{q^{3n+3}}{1-q^{3n+3}}.$$ These correspond to the terms 
(\ref{3.6-0}) to (\ref{3.6-n-2}) in Lemma 3.6.

We complete the proof of Lemma 3.6.

We next consider the {\textbf{Case B}.}

\begin{lemma}\label{Lemma 3.7}%%%%%%%%%%%%%%%%%%%%%%%%%  lemma 3.7
 The generating function for the number of partitions with the given length and alternating sum types $(\Sigma,2)$ and belonging to \textbf{Case B} is the the following $4$ terms.
\begin{equation}\tag{1}\label{3.7-1}
\frac{z^{n+1}q^{3n^{2}+7n+4}}{(zq;q^{3})_{n+1}(q^{3};q^{3})_{n+1}}\frac{n}{z} \, (n\geq 1), \mbox{ for length  } 3n+3,
\end{equation}


\begin{equation}\tag{2}\label{3.7-2}
\frac{z^{n+1}q^{3n^{2}+7n+4}}{(zq;q^{3})_{n+1}(q^{3};q^{3})_{n+1}}\frac{q^{3n+3}}{z}\, (n\geq 0), \mbox { for length } 3n+3,
\end{equation}

\begin{equation}\tag{3}\label{3.7-3}
\frac{z^{n+1}q^{3n^{2}+4n+1}}{(zq;q^{3})_{n+1}(q^{3};q^{3})_{n}} \frac{n}{z} \, (n\geq 1), \mbox{ for length } 3n+1,
\end{equation}


\begin{equation}\tag{4}\label{3.7-4}
\frac{z^{n+1}q^{3n^{2}+4n+1}}{(zq;q^{3})_{n+1}(q^{3};q^{3})_{n}}\frac{q^{3n+3}}{z}  \,(n\geq 0), \mbox{ for length } 3n+2.
\end{equation}
\end{lemma}


\noindent {\it Proof of Lemma 3.7.}

In  \textbf{Case B},  each partition has  only one special unit. We still consider three partitions classes by their length $3n+3$, $3n+2$ and $3n+1$, where $n\geq 0$. In each class, the special unit can be the last basic unit or not.

$\bullet$ {\bf Length $3n+3.$}
$\lambda$ has the form:
$$
\lambda_{1}>\lambda_{2}=\lambda_{3}>\dots\lambda_{3k-3}>{\color{blue}\lambda_{3k-2}\geq \lambda_{3k-1}\overset{2}{>}\lambda_{3k}}\geq\lambda_{3k+1}>\dots>\lambda_{3n+3}, \quad(1\leq k\leq n)
$$
or the form:
$$
\lambda_{1}>\lambda_{2}=\lambda_{3}>\dots\lambda_{3n}>{\color{blue}\lambda_{3n+1}\geq \lambda_{3n+2}\overset{2}{>}\lambda_{3n+3}}, \quad(n \geq 0).
$$
The generating function for the former is
\begin{align*}
\frac{zq}{1-zq}\frac{q^{3}}{1-q^{3}}\dots\frac{q^{3k-3}}{1-q^{3k-3}}&\frac{1}{zq^{3k-2}}q^{6k-2}\frac{1}{1-q^{3k}}\frac{zq^{3k+1}}{1-zq^{3k+1}}\dots\frac{zq^{3n+1}}{1-zq^{3n+1}}\frac{q^{3n+3}}{1-q^{3n+3}},\\
=\frac{zq}{1-zq}\frac{q^{3}}{1-q^{3}}& \frac{zq^{4}}{1-zq^{4}}\dots\frac{zq^{3n+1}}{1-zq^{3n+1}} \frac{q^{3n+3}}{1-q^{3n+3}}\frac{1}{z} \quad(1\leq k \leq n)\\
=&\frac{z^{n+1}q^{3n^{2}+7n+4}}{(zq;q^{3})_{n+1}(q^{3};q^{3})_{n+1}}\frac{1}{z}.
\end{align*}
The special unit can be any $k^{th}$ basic unit, $1\leq k\leq n$,
so we get the term (\ref{3.7-1}) in Lemma 3.7 by adding them.
The generating function for the later is
\begin{align*}
&\frac{zq}{1-zq}\frac{q^{3}}{1-q^{3}}\dots\frac{q^{3n}}{1-q^{3n}}\frac{1}{zq^{3n+1}}q^{6n+4}\frac{q^{3n+3}}{1-q^{3n+3}},\\
&=\frac{zq}{1-zq}\frac{q^{3}}{1-q^{3}} \frac{zq^{4}}{1-zq^{4}}\dots\frac{zq^{3n+1}}{1-zq^{3n+1}} \frac{q^{3n+3}}{1-q^{3n+3}}\frac{q^{3n+3}}{z}\quad (n\geq 0)\\
&=\frac{z^{n+1}q^{3n^{2}+7n+4}}{(zq;q^{3})_{n+1}(q^{3};q^{3})_{n+1}}\frac{q^{3n+3}}{z}.
\end{align*}
We get the term (\ref{3.7-2}) in lemma 3.7. 

$\bullet$ {\bf Length $3n+2.$}
Special unit  in $\lambda$ must be the last basic unit. We note that such partitions can be obtained from partitions belonging to the {\bf Case B} with length $3n+3$  and the special unit being the last basic unit by omitting the part $\lambda_{3n+3}$ and Letting $\lambda_{3n+2}=2$. Hence the generating function for partitions belong to the {\bf Case B} with length $3n+2$ and the special unit being the last basic unit can be obtained from the generating function for partitions belong to {\bf Case B}  with length $3n+3$ and the special is the last basic unit  by omitting the factor $$\frac{q^{3n+3}}{1-q^{3n+3}}.$$  So we get the term (\ref{3.7-4}) in Lemma 3.7.



$\bullet$ {\bf Length $3n+1.$}
In this case, the special unit can not be the last basic unit.  As the analysis above, we find that such partitions can be obtained from partitions belonging to the {\bf Case B} with length $3n+3$  and the special unit being not the last basic unit by omitting the part $\lambda_{3n+3}$ and the part $\lambda_{3n+2}$. Hence the generating function for partitions belong to the {\bf Case B} with length $3n+1$ and the special unit being not the last basic unit can be obtained from the generating function for partitions belong to {\bf Case B}  with length $3n+3$ and the special being not the last basic unit  by omitting the factor $$\frac{q^{3n+3}}{1-q^{3n+3}}.$$
 It corresponds the term (\ref{3.7-3}) in Lemma 3.7.

We complete the proof of Lemma 3.7.


Now the generating function for $a(\Sigma, 2;  n)$ will be the sums of all terms in Lemma 3.4, Lemma 3.5, Lemma 3.6 and Lemma 3.7 involving $3n+3$, $3n+2$, and $3n+1$,  and finally sums them on $n$ over the ranges indicating in these lemmas.
We firstly  compute the sum of terms only involving $\frac{1}{z}$. It is
\begin{align}\nonumber
&\sum_{n\geq 1}\frac{z^{n+1}q^{3n^{2}+7n+4}}{(zq;q^{3})_{n+1}(q^{3};q^{3})_{n+1}}\frac{n}{z}
+\sum_{n\geq 1}\frac{z^{n+1}q^{3n^{2}+4n+1}}{(zq;q^{3})_{n+1}(q^{3};q^{3})_{n}} \frac{n}{z} \\\nonumber
+&\sum_{n\geq 0}\frac{z^{n+1}q^{3n^{2}+7n+4}}{(zq;q^{3})_{n+1}(q^{3};q^{3})_{n+1}}\frac{q^{3n+3}}{z}
+\sum_{n\geq 0}\frac{z^{n+1}q^{3n^{2}+4n+1}}{(zq;q^{3})_{n+1}(q^{3};q^{3})_{n}}\frac{q^{3n+3}}{z}\\\nonumber
+&\sum_{n\geq 2}\frac{z^{n+1}q^{3n^{2}+7n+4}}{(zq;q^{3})_{n+1}(q^{3};q^{3})_{n+1}}\sum_{k=1} ^{n-1}\frac{1-q^{3k}}{zq^{3k}} 
+\sum_{n\geq 1} \frac{z^{n+1}q^{3n^{2}+7n+4}}{(zq;q^{3})_{n+1}(q^{3};q^{3})_{n+1}} \frac{q^{3}(1-q^{3n})}{z} \\\nonumber
+ &\sum_{n\geq 1}\frac{z^{n+1}q^{3n^{2}+4n+1}}{(zq;q^{3})_{n+1}(q^{3};q^{3})_{n}} \frac{q^{3}(1-q^{3n})}{z}
+\sum_{n\geq 2}\frac{z^{n+1}q^{3n^{2}+4n+1}}{(zq;q^{3})_{n+1}(q^{3};q^{3})_{n}}
\sum_{k=1} ^{n-1}\frac{1-q^{3k}}{zq^{3k}}\\\nonumber
=& \sum\mbox{terms}(n\geq 2) + \mbox{ terms }(n=1) + \mbox{terms }(n=0).   \nonumber
\end{align}
After using the identity
$$  
\sum_{n\geq 2}\frac{z^{n+1}q^{3n^{2}+7n+4}}{(zq;q^{3})_{n+1}(q^{3};q^{3})_{n+1}}A
+\sum_{n\geq 2}\frac{z^{n+1}q^{3n^{2}+4n+1}}{(zq;q^{3})_{n+1}(q^{3};q^{3})_{n}} A
=\sum_{n\geq 2}\frac{z^{n+1}q^{3n^{2}+4n+1}}{(zq;q^{3})_{n+1}(q^{3};q^{3})_{n}} \frac{1}{1-q^{3n+3}}A,
$$ we have
\begin{align*}
&\sum\mbox{terms}(n\geq 2)\\\nonumber
&=\sum_{n\geq 2}\frac{z^{n+1}q^{3n^{2}+4n+1}}{(zq;q^{3})_{n+1}(q^{3};q^{3})_{n}} \frac{1}{1-q^{3n+3}}\frac{n}{z}
+\sum_{n\geq 2}\frac{z^{n+1}q^{3n^{2}+4n+1}}{(zq;q^{3})_{n+1}(q^{3};q^{3})_{n}}\frac{1}{1-q^{3n+3}}\frac{q^{3n+3}}{z}\\  \nonumber
&+\sum_{n\geq 2}\frac{z^{n+1}q^{3n^{2}+4n+1}}{(zq;q^{3})_{n+1}(q^{3};q^{3})_{n}} \frac{1}{1-q^{3n+3}} \frac{q^{3}(1-q^{3n})}{z} 
+\sum_{n\geq 2}\frac{z^{n+1}q^{3n^{2}+4n+1}}{(zq;q^{3})_{n+1}(q^{3};q^{3})_{n}} \frac{1}{1-q^{3n+3}}\sum_{k=1} ^{n-1}\frac{1-q^{3k}}{zq^{3k}}\\   \nonumber
&=\sum_{n\geq 2}\frac{z^{n+1}q^{3n^{2}+4n+1}}{(zq;q^{3})_{n+1}(q^{3};q^{3})_{n}}\frac{1}{1-q^{3n+3}}
\left(      \frac{n}{z}   +  \frac{q^{3}(1-q^{3n})}{z}       + \frac{q^{3n+3}}{z}    + \sum_{k=1}^{n-1} \frac{1-q^{3k}}{zq^{3k}}\right)\\\nonumber
&=\sum_{n\geq 2}\frac{zq}{1-zq}\frac{q^{3}}{1-q^{3}} \frac{zq^{4}}{1-zq^{4}}\dots\frac{q^{3n}}{1-q^{3n}}\frac{1}{1-zq^{3n+1}}\\\nonumber
&\left(      \frac{nq^{3n+1}}{1-q^{3n+3}}   +  \frac{q^{3n+4}(1-q^{3n})}{1-q^{3n+3}}       + \frac{q^{6n+4}}{1-q^{3n+3}}    + \frac{q^{4}(1-q^{3n-3})}{(1-q^{3n+3})(1-q^{3})}    +\frac{(1-n)q^{3n+1}}{1-q^{3n+3}}\right)\\ \nonumber
&=\left(\sum_{n\geq 2}\frac{zq}{1-zq}\frac{q^{3}}{1-q^{3}} \frac{zq^{4}}{1-zq^{4}}\dots\frac{q^{3n}}{1-q^{3n}}\frac{1}{1-zq^{3n+1}}\right)\frac{q^{4}}{1-q^{3}}.\\\nonumber
\end{align*} 
And
 \begin{align*}
&\mbox{terms}(n=0) + \mbox{terms}(n=1)\\
=&\frac{1}{1-zq}\frac{q^{4}}{1-q^{3}}+\frac{zq}{1-zq}\frac{q^{3}}{1-q^{3}}\frac{1}{1-zq^{4}}\left(  \frac{q^{10}}{1-q^{6}}+\frac{q^{4}}{1-q^{6}} \right)\\
+&\frac{zq}{1-zq}\frac{q^{3}}{1-q^{3}}\frac{zq^{4}}{1-zq^{4}}\frac{1}{1-q^{6}}\frac{q^{3}(1-q^{3})}{z}\\
=&\frac{1}{1-zq}\frac{q^{4}}{1-q^{3}}+\frac{zq}{1-zq}\frac{q^{3}}{1-q^{3}}\frac{1}{1-zq^{4}}\frac{q^{4}}{1-q^{3}}. 
\end{align*}
Combine it with the last equality above, we get the sum only involving $\frac{1}{z}$ is 
\begin{equation}\label{3.15}%%%%%%%%         (3.15 )          %%%%%
\left(\sum_{n\geq 0}\frac{zq}{1-zq}\frac{q^{3}}{1-q^{3}} \frac{zq^{4}}{1-zq^{4}}\dots\frac{q^{3n}}{1-q^{3n}}\frac{1}{1-zq^{3n+1}}\right)\frac{q^{4}}{1-q^{3}}.
\end{equation}
Now we compute the sum involving $\frac{1}{z^{2}}$ in Lemma 3.4 lemma 3.5 and Lemma 3.6. After using the identity
\begin{align*}
&\left(\frac{z^{n+1}q^{3n^{2}+4n+1}}{(zq;q^{3})_{n+1}(q^{3};q^{3})_{n}} \frac{q^{3n+3}}{1-q^{3n+3}}\right)A+\left(\frac{z^{n+1}q^{3n^{2}+4n+1}}{(zq;q^{3})_{n+1}(q^{3};q^{3})_{n}}\right)A\\
&=\left(\frac{z^{n+1}q^{3n^{2}+4n+1}}{(zq;q^{3})_{n+1}(q^{3};q^{3})_{n}}\frac{1}{1-q^{3n+3}}\right)A.
\end{align*}
We find there are four sums involving $\frac{1}{z^{2}}$, three sums of them are
\begin{align}
&\sum_{n\geq 1}\left(\frac{z^{n+1}q^{3n^{2}+4n+1}}{(zq;q^{3})_{n+1}(q^{3};q^{3})_{n}}\frac{1}{1-q^{3n+3}}\right)\frac{q^{2}}{z^{2}},\\
&\sum_{n\geq 2}\left(\frac{z^{n+1}q^{3n^{2}+4n+1}}{(zq;q^{3})_{n+1}(q^{3};q^{3})_{n}}\frac{1}{1-q^{3n+3}}\right)\sum_{k=1}^{n-1}\frac{1}{z^{2}q^{3k+1}},\\
&\sum_{n\geq 2}\left(\frac{z^{n+1}q^{3n^{2}+4n+1}}{(zq;q^{3})_{n+1}(q^{3};q^{3})_{n}}\frac{1}{1-q^{3n+3}}\right)
\frac{1}{z^{2}}\left( \frac{1}{q} +\frac{1}{q^{4}}\dots\frac{1}{q^{3n-5}}    \right)
\end{align}
and the fourth is the sum of the following $n-2$ terms:
\begin{align}
&\sum_{n\geq 3}\left(\frac{z^{n+1}q^{3n^{2}+4n+1}}{(zq;q^{3})_{n+1}(q^{3};q^{3})_{n}}\frac{1}{1-q^{3n+3}}\right)\sum_{k=1}^{n-2}\frac{1}{z^{2}q^{6k+4}}\label{3.19}\\
+&\sum_{n\geq 4}\left(\frac{z^{n+1}q^{3n^{2}+4n+1}}{(zq;q^{3})_{n+1}(q^{3};q^{3})_{n}}\frac{1}{1-q^{3n+3}}\right)\sum_{k=1}^{n-3}\frac{1}{z^{2}q^{6k+7}}\\
+&\sum_{n\geq 5}\left(\frac{z^{n+1}q^{3n^{2}+4n+1}}{(zq;q^{3})_{n+1}(q^{3};q^{3})_{n}}\frac{1}{1-q^{3n+3}}\right)\sum_{k=1}^{n-4}\frac{1}{z^{2}q^{6k+10}}\\\nonumber
+&\dots\\
+&\sum_{n\geq 3}\left(\frac{z^{n+1}q^{3n^{2}+4n+1}}{(zq;q^{3})_{n+1}(q^{3};q^{3})_{n}}\frac{1}{1-q^{3n+3}}\right)\sum_{k=1}^{n-(n-1)}\frac{1}{z^{2}q^{6k+3n-5}}\label{3.22}
\end{align}
Where the last term (\ref{3.22}) corresponds to the case that the distance between two special units is $n-2$ and the last unit is not a special unit, there is only one case, corresponding to $k=1,$ hence $n>2$. So the range for sum is $n\geq 3$ in (\ref{3.22}). Note that the case of the distance  zero already appears in (\ref{3.19}) (the term of $k=n-2$). Now,
\begin{align*}%%%%%%%     first part
&\mbox{The first sum }=\sum_{n\geq 1}\left(\frac{z^{n+1}q^{3n^{2}+4n+1}}{(zq;q^{3})_{n+1}(q^{3};q^{3})_{n}}\frac{1}{1-q^{3n+3}}\right)\frac{q^{2}}{z^{2}}\\
&=\left(\frac{zq}{1-zq}\frac{q^{3}}{1-q^{3}}\frac{zq^{4}}{1-zq^{4}}\dots\frac{zq^{3n+1}}{1-zq^{3n+1}}\frac{1}{1-q^{3n+3}}    \right)\frac{q^{2}}{z^{2}}\\
&=\left(\frac{zq}{1-zq}\frac{q^{3}}{1-q^{3}}\frac{zq^{4}}{1-zq^{4}}\dots\frac{zq^{3n-2}}{1-zq^{3n-2}}\frac{q^{3n}}{1-q^{3n}}\frac{zq^{3n+1}}{1-zq^{3n+1}}\frac{1}{1-q^{3n+3}}    \right)\frac{q^{2}}{z^{2}}\\
&=\sum_{n\geq 1}\left(\frac{zq}{1-zq}\frac{q^{3}}{1-q^{3}}\dots\frac{q^{3n-3}}{1-q^{3n-3}}\frac{1}{1-zq^{3n-2}}\right)\frac{q^{3n-2}}{1}\cdot\frac{q^{3n}}{1-q^{3n}}\cdot\frac{1}{1-zq^{3n+1}}\cdot\frac{q^{3n+1}}{1-q^{3n+3}}\cdot\frac{q^{2}}{1}\\
&=\sum_{n\geq 1}\left(\frac{zq}{1-zq}\dots\frac{q^{3n-3}}{1-zq^{3n-3}}\frac{1}{1-zq^{3n-2}}\right)\frac{q^{3n-2}}{1}\cdot\frac{q^{3n}}{1-q^{3n}}\cdot(1+zq^{3n+1}+\dots)\cdot\frac{q^{3n+1}}{1-q^{3n+3}}\cdot\frac{q^{2}}{1}\\
&=\sum_{n\geq 1}\left(\frac{zq}{1-zq}\frac{q^{3}}{1-q^{3}}\dots\frac{q^{3n-3}}{1-q^{3n-3}}\frac{1}{1-zq^{3n-2}}\right)\frac{q^{3n-2}}{1}\cdot\frac{q^{3n}}{1-q^{3n}}\cdot1\cdot\frac{q^{3n+1}}{1-q^{3n+3}}\cdot\frac{q^{2}}{1}\\
&+\sum_{n\geq 1}\left(\frac{zq}{1-zq}\frac{q^{3}}{1-q^{3}}\dots\frac{q^{3n-3}}{1-q^{3n-3}}\frac{1}{1-zq^{3n-2}}\right)\frac{q^{3n-2}}{1}\cdot\frac{q^{3n}}{1-q^{3n}}\cdot\frac{zq^{3n+1}}{1-zq^{3n+1}}\cdot\frac{q^{3n+1}}{1-q^{3n+3}}\cdot\frac{q^{2}}{1}\\
&=\sum_{n\geq 1}\left(\frac{zq}{1-zq}\frac{q^{3}}{1-q^{3}}\dots\frac{q^{3n-3}}{1-q^{3n-3}}\frac{1}{1-zq^{3n-2}}\right)\frac{q^{3n-2}}{1}\cdot\frac{q^{3n}}{1-q^{3n}}\cdot1\cdot\frac{q^{3n+1}}{1-q^{3n+3}}\cdot\frac{q^{2}}{1}\\
&+\sum_{n\geq 1}\left(\frac{zq}{1-zq}\frac{q^{3}}{1-q^{3}}\dots\frac{q^{3n}}{1-q^{3n}}\frac{1}{1-zq^{3n+1}}\right)\frac{q^{3n+1}}{1}\cdot\frac{q^{3n+1}}{1-q^{3n+3}}\cdot\frac{q^{2}}{1}\\
&=\sum_{n\geq 1}\left(\frac{zq}{1-zq}\frac{q^{3}}{1-q^{3}}\dots\frac{q^{3n-3}}{1-q^{3n-3}}\frac{1}{1-zq^{3n-2}}\right)\frac{q^{9n+1}}{(1-q^{3n})(1-q^{3n+3})}       \\
&+\sum_{n\geq 1}\left(\frac{zq}{1-zq}\frac{q^{3}}{1-q^{3}}\dots\frac{q^{3n}}{1-q^{3n}}\frac{1}{1-zq^{3n+1}}\right)\frac{q^{6n+4}}{1-q^{3n+3}}. 
\end{align*}


\begin{align*}%%%%%  second part
&\mbox{The second sum }=\\
&\sum_{n\geq 2}\left(\frac{zq}{1-zq}\frac{q^{3}}{1-q^{3}} \frac{zq^{4}}{1-zq^{4}}\dots\frac{zq^{3n+1}}{1-zq^{3n+1}}\frac{1}{1-q^{3n+3}}\right)\sum_{k=1}^{n-1}\frac{1}{z^{2}q^{3k+1}}\\
&=\sum_{n\geq 2}\left(\frac{zq}{1-zq}\dots\frac{q^{3n-3}}{1-q^{3n-3}}\frac{1}{1-zq^{3n-2}}\right)\frac{q^{3n-2}}{1}\cdot\frac{q^{3n}}{1-q^{3n}}\cdot\frac{1}{1-zq^{3n+1}}\cdot\frac{q^{3n+1}}{1-q^{3n+3}}\cdot\sum_{k=1}^{n-1}\frac{1}{q^{3k+1}}\\
&=\sum_{n\geq 2}A_{1}\frac{q^{3n-2}}{1}\cdot\frac{q^{3n}}{1-q^{3n}}\cdot(1+zq^{3n+1}+\dots)\cdot\frac{q^{3n+1}}{1-q^{3n+3}}\cdot\sum_{k=1}^{n-1}\frac{1}{q^{3k+1}}\\
&=\sum_{n\geq 2}A_{1}\frac{q^{3n-2}}{1}\cdot\frac{q^{3n}}{1-q^{3n}}\cdot1\cdot\frac{q^{3n+1}}{1-q^{3n+3}}\cdot\left(\frac{1}{q^{4}}    +\frac{1}{q^{7}}+\dots +\frac{1}{q^{3n-2}}
\right)\\
&+\sum_{n\geq 2}A_{1}\frac{q^{3n-2}}{1}\cdot\frac{q^{3n}}{1-q^{3n}}\cdot\frac{zq^{3n+1}}{1-zq^{3n+1}}\cdot\frac{q^{3n+1}}{1-q^{3n+3}}\cdot\left(\frac{1}{q^{4}}    +\frac{1}{q^{7}}+\dots +\frac{1}{q^{3n-2}}
\right)\\
&=\sum_{n\geq 2}A_{1}\frac{1}{(1-q^{3n})(1-q^{3n+3})}\cdot\left(q^{6n+1}   +q^{6n+4}+\dots +q^{9n-5}
\right)\\
&+\sum_{n\geq 2}A_{2}\frac{1}{1-q^{3n+3}}\cdot\left( q^{3n+4} +q^{3n+7}   +\dots+q^{6n-2}      \right)\\
&=\sum_{n\geq 2}\left(\frac{zq}{1-zq}\frac{q^{3}}{1-q^{3}}\dots\frac{q^{3n-3}}{1-zq^{3n-3}}\frac{1}{1-zq^{3n-2}}\right)\frac{q^{6n+1}(1-q^{3n-3})}{(1-q^{3n})(1-q^{3n+3})(1-q^{3})}    \\
&+\sum_{n\geq 2}\left(\frac{zq}{1-zq}\frac{q^{3}}{1-q^{3}}\dots\frac{q^{3n}}{1-zq^{3n}}\frac{1}{1-zq^{3n+1}}\right)\frac{q^{3n+4}(1-q^{3n-3})}{(1-q^{3n+3})(1-q^{3})}.   \\
&\mbox{ where }A_{1}=\frac{zq}{1-zq}\frac{q^{3}}{1-q^{3}} \dots\frac{q^{3n-3}}{1-q^{3n-3}}\frac{1}{1-zq^{3n-2}},\quad
A_{2}=\frac{zq}{1-zq}\frac{q^{3}}{1-q^{3}}\dots\frac{q^{3n}}{1-q^{3n}}\frac{1}{1-zq^{3n+1}}.
\end{align*}


%%%%%%%%
%%%%%%%  third part



\begin{align*}
&\mbox{The third sum }\\
=&\sum_{n\geq 2}\left(\frac{zq}{1-zq}\frac{q^{3}}{1-q^{3}} \frac{zq^{4}}{1-zq^{4}}\dots\frac{zq^{3n+1}}{1-zq^{3n+1}}\frac{1}{1-q^{3n+3}}\right)
\frac{1}{z^{2}}\left( \frac{1}{q} +\frac{1}{q^{4}}+\dots+\frac{1}{q^{3n-5}}    \right)\\
&=\sum_{n\geq 2}\left(\frac{zq}{1-zq}\dots\frac{1}{1-zq^{3n-2}}\right)\frac{q^{3n-2}}{1}\cdot\frac{q^{3n}}{1-q^{3n}}\cdot\frac{1}{1-zq^{3n+1}}\cdot\frac{q^{3n+1}}{1-q^{3n+3}}\cdot\left( \frac{1}{q} +\frac{1}{q^{4}}+\dots+\frac{1}{q^{3n-5}}    \right)\\
&=\sum_{n\geq 2}A_{1}\frac{q^{3n-2}}{1}\cdot\frac{q^{3n}}{1-q^{3n}}\cdot(1+zq^{3n+1}+\dots)\cdot\frac{q^{3n+1}}{1-q^{3n+3}}\cdot\left( \frac{1}{q} +\frac{1}{q^{4}}+\dots+\frac{1}{q^{3n-5}}    \right)\\
&=\sum_{n\geq 2}A_{1}\frac{q^{3n-2}}{1}\cdot\frac{q^{3n}}{1-q^{3n}}\cdot1\cdot\frac{q^{3n+1}}{1-q^{3n+3}}\cdot\left( \frac{1}{q} +\frac{1}{q^{4}}+\dots+\frac{1}{q^{3n-5}}    \right)\\
&+\sum_{n\geq 2}A_{1}\frac{q^{3n-2}}{1}\cdot\frac{q^{3n}}{1-q^{3n}}\cdot\frac{zq^{3n+1}}{1-zq^{3n+1}}\cdot\frac{q^{3n+1}}{1-q^{3n+3}}\cdot\left( \frac{1}{q} +\frac{1}{q^{4}}+\dots+\frac{1}{q^{3n-5}}    \right)\\
&=\sum_{n\geq 2}A_{1}\frac{1}{(1-q^{3n})(1-q^{3n+3})}\cdot\left(q^{6n+4}   +q^{6n+4}+\dots +q^{9n-2}
\right)\\
&+\sum_{n\geq 2}A_{2}\frac{1}{1-q^{3n+3}}\cdot\left( q^{3n+7} +q^{3n+7}   +\dots+q^{6n+1}      \right)\\
&=\sum_{n\geq 2}\left(\frac{zq}{1-zq}\frac{q^{3}}{1-q^{3}}\dots\frac{q^{3n-3}}{1-zq^{3n-3}}\frac{1}{1-zq^{3n-2}}\right)\frac{q^{6n+4}(1-q^{3n-3})}{(1-q^{3n})(1-q^{3n+3})(1-q^{3})}   \\
&+\sum_{n\geq 2}\left(\frac{zq}{1-zq}\frac{q^{3}}{1-q^{3}}\dots\frac{q^{3n}}{1-zq^{3n}}\frac{1}{1-zq^{3n+1}}\right)\frac{q^{3n+7}(1-q^{3n-3})}{(1-q^{3n+3})(1-q^{3})}.    \\
&\mbox{ where }A_{1}=\frac{zq}{1-zq}\frac{q^{3}}{1-q^{3}} \dots\frac{q^{3n-3}}{1-q^{3n-3}}\frac{1}{1-zq^{3n-2}},\quad
A_{2}=\frac{zq}{1-zq}\frac{q^{3}}{1-q^{3}}\dots\frac{q^{3n}}{1-q^{3n}}\frac{1}{1-zq^{3n+1}}.
\end{align*}




%%%%%%%%%%%%%%%%%%%%%%%%%%        four part
Before computing the fourth sum, we first compute the $d^{th}$ term in the fourth sum:
\begin{align*}
&\mbox{The $d^{th}$ term}\\
&=\sum_{n\geq d+2}\left(\frac{zq}{1-zq}\frac{q^{3}}{1-q^{3}} \frac{zq^{4}}{1-zq^{4}}\dots\frac{zq^{3n+1}}{1-zq^{3n+1}}\frac{1}{1-q^{3n+3}}\right)
\sum_{k=1}^{n-d-1}\frac{1}{z^{2}}\frac{1}{q^{6k+3d+1}}\\
&=\sum_{n\geq d+2}\left(\frac{zq}{1-zq}\dots\frac{1}{1-zq^{3n-2}}\right)\frac{q^{3n-2}}{1}\cdot\frac{q^{3n}}{1-q^{3n}}\cdot\frac{1}{1-zq^{3n+1}}\cdot\frac{q^{3n+1}}{1-q^{3n+3}}\cdot\sum_{k=1}^{n-d-1}\frac{1}{q^{6k+3d+1}}\\
&=\sum_{n\geq d+2}A_{1}\frac{q^{3n-2}}{1}\cdot\frac{q^{3n}}{1-q^{3n}}\cdot(1+zq^{3n+1}+\dots)\cdot\frac{q^{3n+1}}{1-q^{3n+3}}\cdot\sum_{k=1}^{n-d-1}\frac{1}{q^{6k+3d+1}}\\
&=\sum_{n\geq d+2}A_{1}\frac{q^{3n-2}}{1}\cdot\frac{q^{3n}}{1-q^{3n}}\cdot1\cdot\frac{q^{3n+1}}{1-q^{3n+3}}\cdot\sum_{k=1}^{n-d-1}\frac{1}{q^{6k+3d+1}}\\
&+\sum_{n\geq d+2}A_{1}\frac{q^{3n-2}}{1}\cdot\frac{q^{3n}}{1-q^{3n}}\cdot\frac{zq^{3n+1}}{1-zq^{3n+1}}\cdot\frac{q^{3n+1}}{1-q^{3n+3}}\cdot\sum_{k=1}^{n-d-1}\frac{1}{q^{6k+3d+1}}\\
&=\sum_{n\geq d+2}A_{1}\frac{1}{(1-q^{3n})(1-q^{3n+3})}\cdot\left(q^{3n+3d+4}   +q^{3n+3d+10}+\dots +q^{9n-3d-8}
\right)\\
&+\sum_{n\geq d+2}A_{2}\frac{1}{1-q^{3n+3}}\cdot\left( q^{3d+7} +q^{3d+13}   +\dots+q^{6n-3d-5}      \right)\\
&=\sum_{n\geq d+2}\left(\frac{zq}{1-zq}\frac{q^{3}}{1-q^{3}}\dots\frac{q^{3n-3}}{1-zq^{3n-3}}\frac{1}{1-zq^{3n-2}}\right)\frac{q^{3n+3d+4}(1-q^{6n-6d-6})}{(1-q^{3n})(1-q^{3n+3})(1-q^{6})}\\
&+\sum_{n\geq d+2}\left(\frac{zq}{1-zq}\frac{q^{3}}{1-q^{3}}\dots\frac{q^{3n}}{1-zq^{3n}}\frac{1}{1-zq^{3n+1}}\right)\frac{q^{3d+7}(1-q^{6n-6d-6})}{(1-q^{3n+3})(1-q^{6})}.\\
&\mbox{ where }A_{1}=\frac{zq}{1-zq}\frac{q^{3}}{1-q^{3}} \dots\frac{q^{3n-3}}{1-q^{3n-3}}\frac{1}{1-zq^{3n-2}},\quad
A_{2}=\frac{zq}{1-zq}\frac{q^{3}}{1-q^{3}}\dots\frac{q^{3n}}{1-q^{3n}}\frac{1}{1-zq^{3n+1}}.
\end{align*}
Therefore, the fourth sum $=$ 
\begin{align*}
&\sum_{n\geq 3}\left(\frac{zq}{1-zq}\frac{q^{3}}{1-q^{3}}\dots\frac{q^{3n-3}}{1-zq^{3n-3}}\frac{1}{1-zq^{3n-2}}\right)\sum_{d=1}^{n-2}\frac{q^{3n+3d+4}(1-q^{6n-6d-6})}{(1-q^{3n})(1-q^{3n+3})(1-q^{6})}\\
&+\sum_{n\geq 3}\left(\frac{zq}{1-zq}\frac{q^{3}}{1-q^{3}}\dots\frac{q^{3n}}{1-zq^{3n}}\frac{1}{1-zq^{3n+1}}\right)\sum_{d=1}^{n-2}\frac{q^{3d+7}(1-q^{6n-6d-6})}{(1-q^{3n+3})(1-q^{6})}\\
&=\sum_{n\geq 3}\left(\frac{zq}{1-zq}\frac{q^{3}}{1-q^{3}}\dots\frac{q^{3n-3}}{1-zq^{3n-3}}\frac{1}{1-zq^{3n-2}}\right)\frac{q^{3n+7}(1-q^{3n-3})(1-q^{3n-6})}{(1-q^{3n})(1-q^{3n+3})(1-q^{6})(1-q^{3})}  \\
&+\sum_{n\geq 3}\left(\frac{zq}{1-zq}\frac{q^{3}}{1-q^{3}}\dots\frac{q^{3n}}{1-zq^{3n}}\frac{1}{1-zq^{3n+1}}\right)\frac{q^{10}(1-q^{3n-3})(1-q^{3n-6})}{(1-q^{3n+3})(1-q^{6})(1-q^{3})}.
\end{align*}
Now the sum of four sums is 
$
\mbox{the first sum} + \mbox{the second sum} +\mbox{ the third sum} +\mbox{the fourth sum}
$ indicated above. Since in the four sums,  the common range of $n$ is $n\geq 3$, so we first consider the four sums over all  $n\geq 3$, and then consider terms corresponding to $n=1$ and $n=2$.
\begin{align*}
&\mbox{The first sum restrict to } n\geq 3 + \mbox{the second sum restrict to }n\geq 3\\
  +&\mbox{ the third sum restrict to } n\geq3+\mbox{the fourth sum restrict to }n\geq 3\\
&=\sum_{n\geq 3}\left(\frac{zq}{1-zq}\frac{q^{3}}{1-q^{3}}\dots\frac{q^{3n-3}}{1-zq^{3n-3}}\frac{1}{1-zq^{3n-2}}\right)\frac{q^{9n+1}}{(1-q^{3n})(1-q^{3n+3})}     \\
&+\sum_{n\geq 3}\left(\frac{zq}{1-zq}\frac{q^{3}}{1-q^{3}}\dots\frac{q^{3n}}{1-zq^{3n}}\frac{1}{1-zq^{3n+1}}\right)\frac{q^{6n+4}}{1-q^{3n+3}}\\
&+\sum_{n\geq 3}\left(\frac{zq}{1-zq}\frac{q^{3}}{1-q^{3}}\dots\frac{q^{3n-3}}{1-zq^{3n-3}}\frac{1}{1-zq^{3n-2}}\right)\frac{q^{6n+1}(1-q^{3n-3})}{(1-q^{3n})(1-q^{3n+3})(1-q^{3})}  \\
&+\sum_{n\geq 3}\left(\frac{zq}{1-zq}\frac{q^{3}}{1-q^{3}}\dots\frac{q^{3n}}{1-zq^{3n}}\frac{1}{1-zq^{3n+1}}\right)\frac{q^{3n+4}(1-q^{3n-3})}{(1-q^{3n+3})(1-q^{3})}  \\
&+\sum_{n\geq 3}\left(\frac{zq}{1-zq}\frac{q^{3}}{1-q^{3}}\dots\frac{q^{3n-3}}{1-zq^{3n-3}}\frac{1}{1-zq^{3n-2}}\right)\frac{q^{6n+4}(1-q^{3n-3})}{(1-q^{3n})(1-q^{3n+3})(1-q^{3})}  \\
&+\sum_{n\geq 3}\left(\frac{zq}{1-zq}\frac{q^{3}}{1-q^{3}}\dots\frac{q^{3n}}{1-zq^{3n}}\frac{1}{1-zq^{3n+1}}\right)\frac{q^{3n+7}(1-q^{3n-3})}{(1-q^{3n+3})(1-q^{3})}.\\
&+\sum_{n\geq 3}\left(\frac{zq}{1-zq}\frac{q^{3}}{1-q^{3}}\dots\frac{q^{3n-3}}{1-zq^{3n-3}}\frac{1}{1-zq^{3n-2}}\right)\frac{q^{3n+7}(1-q^{3n-3})(1-q^{3n-6})}{(1-q^{3n})(1-q^{3n+3})(1-q^{6})(1-q^{3})} \\
&+\sum_{n\geq 3}\left(\frac{zq}{1-zq}\frac{q^{3}}{1-q^{3}}\dots\frac{q^{3n}}{1-zq^{3n}}\frac{1}{1-zq^{3n+1}}\right)\frac{q^{10}(1-q^{3n-3})(1-q^{3n-6})}{(1-q^{3n+3})(1-q^{6})(1-q^{3})}\\
&=\sum_{n\geq 3}A_{1}\left(\frac{q^{9n+1}}{(1-q^{3n})(1-q^{3n+3})}\right.+\frac{q^{6n+1}(1-q^{3n-3})}{(1-q^{3n})(1-q^{3n+3})(1-q^{3})}\\
&+\frac{q^{6n+4}(1-q^{3n-3})}{(1-q^{3n})(1-q^{3n+3})(1-q^{3})}+\left.\frac{q^{3n+7}(1-q^{3n-3})(1-q^{3n-6})}{(1-q^{3})(1-q^{6})(1-q^{3n})(1-q^{3n+3})}\right) \\
&+\sum_{n\geq 3}A_{2}\left(\frac{q^{6n+4}}{1-q^{3n+3}}\right.+\frac{q^{3n+4}(1-q^{3n-3})}{(1-q^{3n+3})(1-q^{3})}\\
&+\frac{q^{3n+7}(1-q^{3n-3})}{(1-q^{3n+3})(1-q^{3})}+\left.\frac{q^{10}(1-q^{3n-3})(1-q^{3n-6})}{(1-q^{3})(1-q^{6})(1-q^{3n+3})}\right) \\
&=\sum_{n\geq 3}A_{1}\frac{q^{3n+7}}{(1-q^{3})(1-q^{6})}+\sum_{n\geq 3}A_{2}\frac{q^{10}(1-q^{3n})}{(1-q^{3})(1-q^{6})}.\\
&\mbox{ where }A_{1}=\frac{zq}{1-zq}\frac{q^{3}}{1-q^{3}} \dots\frac{q^{3n-3}}{1-q^{3n-3}}\frac{1}{1-zq^{3n-2}},\quad
A_{2}=\frac{zq}{1-zq}\frac{q^{3}}{1-q^{3}}\dots\frac{q^{3n}}{1-q^{3n}}\frac{1}{1-zq^{3n+1}}.
 \end{align*}
But 
$$
\sum_{n\geq 3}A_{1}\frac{q^{3n+7}}{(1-q^{3})(1-q^{6})}=\sum_{n\geq 2}A_{2}\frac{q^{3n+10}}{(1-q^{3})(1-q^{6})},
$$
hence 
\begin{align}
&\mbox{The first sum restrict to } n\geq 3 + \mbox{the second sum restrict to }n\geq 3\nonumber\\
  +&\mbox{ the third sum restrict to } n\geq3+\mbox{the fourth sum restrict to }n\geq 3\nonumber\\
 &= \frac{zq}{1-zq}\frac{q^{3}}{1-q^{3}}\frac{zq^{4}}{1-zq^{4}}\frac{q^{6}}{1-q^{6}}\frac{1}{1-zq^{7}}\frac{q^{16}}{(1-q^{3})(1-q^{6})} \nonumber\\
 &+\sum_{n\geq 3}A_{2}\frac{q^{3n+10}}{(1-q^{3})(1-q^{6})} +\sum_{n\geq 3}A_{2}\frac{q^{10}(1-q^{3n})}{(1-q^{3})(1-q^{6})} \nonumber\\
 &= \frac{zq}{1-zq}\frac{q^{3}}{1-q^{3}}\frac{zq^{4}}{1-zq^{4}}\frac{q^{6}}{1-q^{6}}\frac{1}{1-zq^{7}}\frac{q^{16}}{(1-q^{3})(1-q^{6})}+\sum_{n\geq 3}A_{2}\frac{q^{10}}{(1-q^{3})(1-q^{6})}.\label{3.23}
 \end{align}
 Now we consider the sum of  terms $n\leq 2$, their sum is 
 \begin{align*}
&\mbox{The first sum restrict to } n<3 + \mbox{the second sum restrict to }n<3\\
  +&\mbox{ the third sum restrict to } n<3+\mbox{the fourth sum restrict to }n<3\\
&=\sum_{n=1}^{2}\left(\frac{zq}{1-zq}\frac{q^{3}}{1-q^{3}}\dots\frac{q^{3n-3}}{1-zq^{3n-3}}\frac{1}{1-zq^{3n-2}}\right)\frac{q^{9n+1}}{(1-q^{3n})(1-q^{3n+3})}     \\
&+\sum_{n=1}^{2}\left(\frac{zq}{1-zq}\frac{q^{3}}{1-q^{3}}\dots\frac{q^{3n}}{1-zq^{3n}}\frac{1}{1-zq^{3n+1}}\right)\frac{q^{6n+4}}{1-q^{3n+3}}\\
&+\sum_{n=2}^{2}\left(\frac{zq}{1-zq}\frac{q^{3}}{1-q^{3}}\dots\frac{q^{3n-3}}{1-zq^{3n-3}}\frac{1}{1-zq^{3n-2}}\right)\frac{q^{6n+1}(1-q^{3n-3})}{(1-q^{3n})(1-q^{3n+3})(1-q^{3})}  \\
&+\sum_{n=2}^{2}\left(\frac{zq}{1-zq}\frac{q^{3}}{1-q^{3}}\dots\frac{q^{3n}}{1-zq^{3n}}\frac{1}{1-zq^{3n+1}}\right)\frac{q^{3n+4}(1-q^{3n-3})}{(1-q^{3n+3})(1-q^{3})}  \\
&+\sum_{n=2}^{2}\left(\frac{zq}{1-zq}\frac{q^{3}}{1-q^{3}}\dots\frac{q^{3n-3}}{1-zq^{3n-3}}\frac{1}{1-zq^{3n-2}}\right)\frac{q^{6n+4}(1-q^{3n-3})}{(1-q^{3n})(1-q^{3n+3})(1-q^{3})}  \\
&+\sum_{n=2}^{2}\left(\frac{zq}{1-zq}\frac{q^{3}}{1-q^{3}}\dots\frac{q^{3n}}{1-zq^{3n}}\frac{1}{1-zq^{3n+1}}\right)\frac{q^{3n+7}(1-q^{3n-3})}{(1-q^{3n+3})(1-q^{3})}\\
&+\frac{zq}{1-zq}\frac{q^{3}}{1-q^{3}}\frac{zq^{4}}{1-zq^{4}}\frac{q^{6}}{1-q^{6}}\frac{1}{1-zq^{7}}\frac{q^{16}}{(1-q^{3})(1-q^{6})}
 \end{align*}
(Note that last term above appears in (\ref{3.23}). )
\begin{align}
&=\frac{zq}{1-zq}\frac{q^{10}}{(1-q^{3})(1-q^{6})} +\frac{zq}{1-zq}\frac{q^{3}}{1-q^{3}}\frac{1}{1-zq^{4}}\frac{q^{19}}{(1-q^{6})(1-q^{9})} +\frac{zq}{1-zq}\frac{q^{3}}{1-q^{3}}\frac{1}{1-zq^{4}}\frac{q^{10}}{1-q^{6}} \nonumber\\
&+\frac{zq}{1-zq}\frac{q^{3}}{1-q^{3}}\frac{1}{1-zq^{4}}\frac{q^{13}}{(1-q^{6})(1-q^{9})}
+\frac{zq}{1-zq}\frac{q^{3}}{1-q^{3}}\frac{1}{1-zq^{4}}\frac{q^{6}}{1-q^{6}}\frac{1}{1-zq^{7}}\frac{q^{10}}{1-q^{9}}\nonumber\\
&+\frac{zq}{1-zq}\frac{q^{3}}{1-q^{3}}\frac{1}{1-zq^{4}}\frac{q^{16}}{(1-q^{6})(1-q^{9})} 
+\frac{zq}{1-zq}\frac{q^{3}}{1-q^{3}}\frac{1}{1-zq^{4}}\frac{q^{6}}{1-q^{6}}\frac{1}{1-zq^{7}}\frac{q^{13}}{1-q^{9}}\nonumber\\
&+\frac{zq}{1-zq}\frac{q^{3}}{1-q^{3}}\frac{zq^{4}}{1-zq^{4}}\frac{q^{6}}{1-q^{6}}\frac{1}{1-zq^{7}}\frac{q^{16}}{(1-q^{3})(1-q^{6})}\nonumber\\
&=\left(\frac{zq}{1-zq} +\frac{zq}{1-zq}\frac{q^{3}}{1-q^{3}}\frac{1}{1-zq^{4}}+\frac{zq}{1-zq}\frac{q^{3}}{1-q^{3}}\frac{1}{1-zq^{4}}\frac{q^{6}}{1-q^{6}}\frac{1}{1-zq^{7}}  \right)\frac{q^{10}}{(1-q^{3})(1-q^{6})}.\label{3.24}
\end{align} 
Combine (\ref{3.23}) with (\ref{3.24}), we get the sum only involving $\frac{1}{z^{2}}$ is
\begin{equation}\label{3.25}
\left(\sum_{n\geq 0}\frac{zq}{1-zq}\frac{q^{3}}{1-q^{3}} \frac{zq^{4}}{1-zq^{4}}\dots\frac{q^{3n}}{1-q^{3n}}\frac{1}{1-zq^{3n+1}}\right)\frac{q^{10}}{(1-q^{3})(1-q^{6})}.
\end{equation}
Therefore, the generating function for the partitions with repeated times $\leq2$ and alternating sum types $(\Sigma,2)$ is
\begin{align*}
&\sum_{\Sigma, n\geq 0}a(\Sigma, 2;n)z^{\Sigma}q^{n}
=\mbox{the sums involving} \frac{1}{z} + \mbox{ the sums involving} \frac{1}{z^{2}}=(\ref{3.15})+(\ref{3.25})\\
&=\left(\sum_{n\geq 0}\frac{zq}{1-zq}\frac{q^{3}}{1-q^{3}} \frac{zq^{4}}{1-zq^{4}}\dots\frac{q^{3n}}{1-q^{3n}}\frac{1}{1-zq^{3n+1}}\right)\frac{q^{4}}{(1-q^{3})(1-q^{6})}.
\end{align*}
But the function
$$
\sum_{n\geq 0}\frac{zq}{1-zq}\frac{q^{3}}{1-q^{3}} \frac{zq^{4}}{1-zq^{4}}\dots\frac{q^{3n}}{1-q^{3n}}\frac{1}{1-zq^{3n+1}}
$$
is exactly the generating function of $a(\Sigma, 0; n)$!  Since the general term 
$$
\frac{zq}{1-zq}\frac{q^{3}}{1-q^{3}} \frac{zq^{4}}{1-zq^{4}}\dots\frac{q^{3n}}{1-q^{3n}}\frac{1}{1-zq^{3n+1}}
$$
generates all partitions with pure type $(\Sigma, 0)$  with length $3n$ or $3n+1$:  the generating function for partitions with pure type $(\Sigma, 0)$ and length $3n$ (by considering their conjugates) corresponds to 
$$
\frac{zq}{1-zq}\frac{q^{3}}{1-q^{3}} \frac{zq^{4}}{1-zq^{4}}\dots\frac{q^{3n}}{1-q^{3n}}
$$
and the generating function for partitions with pure type $(\Sigma, 0)$ and length $3n+1$ corresponds to
$$
\frac{zq}{1-zq}\frac{q^{3}}{1-q^{3}} \frac{zq^{4}}{1-zq^{4}}\dots\frac{q^{3n}}{1-q^{3n}}(zq^{3n+1}+z^{2}q^{6n+2}+\dots).$$
Hence we proved $a(\Sigma, 2;n)$ satisfies the same relation as $b(\Sigma, 2;n)$.   

The proof of Theorem 3.1 is completed.

\noindent {\it Proof of Theorem 3.2.}

We prove Theorem 3.2 for alternating sum types $(\Sigma, 1,0, \dots, 0)$. Proofs for other cases are similar.
A basic unit has length $m$ consisting in two types:
$$
\lambda_{1} \geq \lambda_{2}\overset{1}{>}\lambda_{3}=\dots=\lambda_{m} \mbox{  or } \lambda_{1}>\lambda_{2}=\lambda_{3}=\dots=\lambda_{m}.
$$
We call the first type is a special unit. Note that there is only one special unit in each partition if its  alternating sum type is $(\Sigma, 1,0, \dots, 0)$. The alternating sum types  $(\Sigma, 1,0, \dots, 0)$ above imply that all partitions can be classified into $3$ classes according to their length $mn+1$,  $mn+2$ and $mn+m$, $n\geq 0$.

$\bullet$ {\bf Length $mn+m.$} Such a partition has $n+1$ basic units.
If the special unit is the $k^{th}$ unit but not the last basic unit, then such it has the form:
$$
\dots>
{\color{blue}\lambda_{(k-1)m+1}\geq \lambda_{(k-1)m+2}\overset{1}{>}\lambda_{(k-1)m+3}=\dots=\lambda_{km}}\geq \lambda_{km+1}>\dots \quad    \quad(1\leq k\leq n)
$$
or the form:
$$
\dots=\lambda_{nm}>
{\color{blue}\lambda_{nm+1}\geq \lambda_{nm+2}\overset{1}{>}\lambda_{nm+3}=\dots=\lambda_{nm+n}},
$$
if the special unit is the last basic unit.

The generating function for the former is
\begin{align*}
&\frac{zq}{1-zq}\frac{q^{m}}{1-q^{m}}\frac{zq^{m+1}}{1-zq^{m+1}}\frac{q^{2m}}{1-q^{2m}}\dots\frac{q^{(k-1)m}}{1-q^{(k-1)m}}\\\nonumber\\
\times&\frac{1}{1-zq^{(k-1)m+1}}q^{(k-1)m+2}\frac{1}{1-q^{km}}\frac{zq^{km+1}}{1-zq^{km+1}}\dots \frac{q^{nm+m}}{1-q^{nm+m}}\\ \nonumber\\
=&\frac{zq}{1-zq}\dots\frac{q^{mn+m}}{1-q^{mn+m}}\left(  \frac{1}{zq^{(k-1)m+1}} \cdot q^{(k-1)m+2} \cdot \frac{1}{q^{km}}  \right)\nonumber\\
=&\frac{zq}{1-zq}\dots\frac{q^{mn+m}}{1-q^{mn+m}}\frac{1}{zq^{km-1}},\label{3.26} \quad(n\geq1),
\end{align*}
and the generating function for the later is
\begin{align*}
&\frac{zq}{1-zq}\frac{q^{m}}{1-q^{m}}\frac{zq^{m+1}}{1-zq^{m+1}}\frac{q^{2m}}{1-q^{2m}}\dots\frac{q^{nm}}{1-q^{nm}}\frac{1}{1-zq^{nm+1}}q^{nm+2}\frac{q^{nm+m}}{1-q^{nm+m}}\nonumber\\
&=\frac{zq}{1-zq}\dots\frac{q^{mn+m}}{1-q^{mn+m}}\frac{q}{z},  \quad (n\geq 0).
\end{align*}
Hence the generating function for partitions with length $mn+m$ is
\begin{equation}
\sum_{k=1}^{n}\frac{zq}{1-zq}\dots\frac{q^{mn+m}}{1-q^{mn+m}}\frac{1}{zq^{km-1}}(n\geq 1)+\frac{zq}{1-zq}\dots\frac{q^{mn+m}}{1-q^{mn+m}}\frac{q}{z}(n\geq0).\label{3.26}
\end{equation}



$\bullet$ {\bf Length $mn+2$.}
In this case, the special unit must be the last unit. and $\lambda_{mn+2}=1$. The generating function for such partitions is
\begin{align}
\frac{zq}{1-zq}\dots\frac{q^{mn}}{1-q^{mn}}\cdot\frac{1}{1-zq^{mn+1}}\cdot q^{mn+2}(n\geq0).\label{3.27}
\end{align}




$\bullet$ {\bf Length $mn+1.$}
In this case, the special unit can not be the last unit, we can obtained such partition from partitions with length $mn+m$ and the special unit being not the last unit by omitting the parts from $\lambda_{mn+2}$ to $\lambda_{mn+m}$. Therefore, the generating function for partitions with length  $nm+1$ and the special unit being not the last unit can be obtained from the generating function for partitions with length $nm+m$ and the special unit being not the last unit by omitting the factor
$$
\frac{q^{nm+m}}{1-q^{nm+m}}.
$$
It will be
\begin{equation}
\sum_{k=1}^{n}\frac{zq}{1-zq}\dots\frac{zq^{mn+1}}{1-zq^{mn+1}}\frac{1}{zq^{km-1}},    \quad(n\geq 1).\label{3.28}
\end{equation}
Combining (\ref{3.26}), (\ref{3.27}) and (\ref{3.28}) and summing them in the range $n\geq 0$, we get the generating function for $a(\Sigma, 1, 0\dots, 0;  n)$.  It is the sum of

\begin{align*}\nonumber
&\sum_{n\geq1}\left(\sum_{k=1}^{n}\frac{zq}{1-zq}\dots\frac{q^{mn+m}}{1-q^{mn+m}}\frac{1}{zq^{km-1}}+\frac{zq}{1-zq}\dots\frac{q^{mn+m}}{1-q^{mn+m}}\frac{q}{z}\right.\\
+&\frac{zq}{1-zq}\dots\frac{q^{mn}}{1-q^{mn}}\cdot\frac{1}{1-zq^{mn+1}}\cdot q^{mn+2}
+\sum_{n\geq1}^{}\frac{zq}{1-zq}\dots\frac{zq^{mn+1}}{1-zq^{mn+1}}\left.\frac{1}{zq^{km-1}}\right)\\
=&\sum_{n\geq1}^{}\frac{zq}{1-zq}\dots\frac{zq^{mn+1}}{1-zq^{mn+1}}\left( \sum_{k=1}^{n}\frac{1}{zq^{km-1}}+\frac{q^{mn+2}}{zq^{mn+1}} +\frac{q^{mn+m}}{1-q^{mn+m}}\frac{q}{z}  +\frac{q^{mn+m}}{1-q^{mn+m}}\sum_{k=1}^{n}\frac{1}{zq^{km-1}}\right)\\
=&\sum_{n\geq1}^{}\frac{zq}{1-zq}\dots\frac{zq^{mn+1}}{1-zq^{mn+1}}\left( \frac{1}{1-q^{mn+m}}\sum_{k=1} ^{n} \frac{1}{zq^{km-1}} +\frac{1}{1-q^{mn+m}} \frac{q}{z} \right)\\
=&\sum_{n\geq1}^{}\frac{zq}{1-zq}\dots\frac{zq^{mn+1}}{1-zq^{mn+1}}\cdot \frac{1}{z(1-q^{mn+m})}\cdot\left(   \sum_{k=1}^{n} \frac{1}{q^{km-1}} +q \right)\\
=&\sum_{n\geq1}^{}\frac{zq}{1-zq}\dots\frac{1}{1-zq^{mn+1}}\cdot\frac{q^{mn+1}}{1-q^{mn+m}}\cdot\left(      \sum_{k=1}^{n} \frac{1}{q^{km-1}} +q  \right)\\
=&\sum_{n\geq1}^{}\frac{zq}{1-zq}\dots\frac{1}{1-zq^{mn+1}}\cdot\frac{1}{1-q^{mn+m}}\cdot\left( q^{2} +q^{m+2}+\dots+q^{mn-m+2}+q^{mn+2} \right)\\
=&\sum_{n\geq1}^{}\frac{zq}{1-zq}\dots\frac{1}{1-zq^{mn+1}}\cdot\frac{1}{1-q^{mn+m}}\cdot q^{2}\cdot \frac{1-q^{mn+m}}{1-q^{m}} \\
=&\sum_{n\geq1}^{}\frac{zq}{1-zq}\dots\frac{1}{1-zq^{mn+1}}\cdot \frac{q^{2}}{1-q^{m}},
\end{align*}
and the term corresponding to $n=0$, which is
\begin{align*}
\frac{1}{1-zq}\cdot q^{2}+\frac{zq}{1-zq}\frac{q^{m}}{1-q^{m}}\cdot \frac{q}{z}
=\frac{1}{1-zq}\cdot \frac{q^{2}}{1-q^{m}}.
\end{align*}
Hence, the generating function for $a(\Sigma, 1, 0,\dots,0;    n)$ is
$$
\sum_{n\geq0}^{}\frac{zq}{1-zq}\dots\frac{1}{1-zq^{mn+1}}\cdot \frac{q^{2}}{1-q^{m}}.$$
Comparing with (\ref{3.3}) for $i=1, j=2$, it shows that Theorem 3.2 is true,  since
$$
\sum_{n\geq0}^{}\frac{zq}{1-zq}\dots\frac{1}{1-zq^{mn+1}}
$$
generates all partitions with alternating sum types $(\Sigma, 0, 0,\dots,0)$. 

The proof of Theorem 3.2 is completed.




\section{$q$-difference equations for moduli three }
Many theorems on partitions are proved by $q$-difference equations method. In this section, we use this method to study our conjecture  for moduli $m=3$, we can get three $q$-difference equations related to the generating function for partitions in $P$. 

Let $a_{N}(\Sigma_{1},  \Sigma_{2};  n)$ denote the number of partitions of $n$ with alternating sum types $(\Sigma_{1}, \Sigma_{2})$ and the length $N$,  let $a_{N,k}(\Sigma_{1}, \Sigma_{2}; n)$ denote the number of partitions of $n$ with alternating sum types $(\Sigma_{1}, \Sigma_{2})$ and length  $N$ and the last part $\lambda_{N}=k$, $k\geq 1$. Note that $a_{N,k}(\Sigma_{1}, \Sigma_{2}; n)=0$ , if $k>n$ or $n<0$, also $a_{N,k}(\Sigma_{1}, \Sigma_{2}; n)=0$ , if $\Sigma_{1}<0$ or $\Sigma_{2}<0$. 
\begin{lemma}
For all $\Sigma_{1}\geq 0$, $\Sigma_{2}\geq 0$ and $n\geq 1$,
\begin{align}
a_{3N}(\Sigma_{1}, \Sigma_{2}; n)&=\sum_{k\geq 1}a_{3N-1}(\Sigma_{1}, \Sigma_{2}; n-3Nk) +\sum_{k\geq 1}a_{3N-2}(\Sigma_{1}, \Sigma_{2}; n-3Nk),\label{4.1}\\
a_{3N+1}(\Sigma_{1}, \Sigma_{2}; n)&=\sum_{k\geq 1}a_{3N}(\Sigma_{1}-k, \Sigma_{2}; n-3Nk-k) +\sum_{k\geq 1}a_{3N-1}(\Sigma_{1}-k, \Sigma_{2}; n-3Nk-k),\label{4.2}\\
a_{3N+2}(\Sigma_{1}, \Sigma_{2}; n)&=\sum_{k\geq 1}a_{3N+1}(\Sigma_{1}, \Sigma_{2}-k; n-3Nk-2k) +\sum_{k\geq 1}a_{3N}(\Sigma_{1}, \Sigma_{2}-k; n-3Nk-2k).\label{4.3}
\end{align}
\end{lemma}

\noindent {\it Proof.}
 Let us firstly prove (\ref{4.1}). The set of partitions counted by $a_{3N}(\Sigma_{1}, \Sigma_{2}; n)$ can be classified according to the size of the last part $\lambda_{N}$ and we have $a_{3N}(\Sigma_{1}, \Sigma_{2}; n)=\sum_{k\geq 1} a_{3N, k}(\Sigma_{1},\Sigma_{2}; n)$. Let us consider a partition counted by $a_{3N,k}(\Sigma_{1},\Sigma_{2};  n)$, note that $\lambda_{3N}=k$.   By the condition each part repeated at most two times, we have two cases:

$\bullet$  $\lambda_{3N-2}\geq \lambda_{3N-1}>\lambda_{3N}.$  We subtract $k$ from each part of this partition, we get a partition of length $3N-1$ and each part still repeated at most two times. Moreover, subtracting $k$ from each part does mot change its alternating sum type,  hence we get a partition counted by $a_{3N-1}(\Sigma_{1}, \Sigma_{2}; n-3Nk)$.
Conversely, we can add $k$ to each part of a partition counted by $a_{3N-1}(\Sigma_{1},\Sigma_{2}; n-3Nk)$ and next add a new part $\lambda_{3N}=k$, so we get a partition counted by $a_{3N,k}(\Sigma_{1}, \Sigma_{2}; n)$. 

$\bullet$ $\lambda_{3N-2}> \lambda_{3N-1}=\lambda_{3N}.$   As before, we  still subtract $k$ from each part of this partition, however, we get a partition of $n-3Nk$ with length $3N-2$ and each part repeated at most times, moreover, its alternating sum type is still $(\Sigma_{1},\Sigma_{2})$, hence  a partition counted by $a_{3n-2}(\Sigma_{1},\Sigma_{2}; n-3Nk)$. Conversely, we can get a partition counted by 
$a_{3N,k}(\Sigma_{1},\Sigma_{2}; n)$ by adding $k$ to each part of a partition counted by $a_{3N-2}(\Sigma_{1},\Sigma_{2}; n-3Nk)$ and taking $\lambda_{3N-1}=\lambda_{3N}=k$. 

Combine with these two cases, we have
\begin{align*}
a_{3N}(\Sigma_{1},\Sigma_{2};n)
&=\sum_{k\geq 1}a_{3N,k}(\Sigma_{1}, \Sigma_{2};n)\\
&=\sum_{k\geq 1}\bigg( a_{3N-1}(\Sigma_{1},\Sigma_{2}; n-3Nk)   + a_{3N-2} (\Sigma_{1},\Sigma_{2};n-3Nk)           \bigg),
\end{align*}
which is (\ref{4.1}).

The proofs of equations (\ref{4.2}) and (\ref{4.3}) are similar to the proof of (\ref{4.1}), except two things will change. One is that when we subtract $k$ from each part, the alternating sum type will change. In case of (\ref{4.2}), the length is $3N+1$, the alternating sum type will change to $(\Sigma_{1}-k,\Sigma_{2})$, if $\lambda_{3N+1}=k$. In case of (\ref{4.3}), the length is $3N+2$, its alternating sum type will change to  $(\Sigma_{1},  \Sigma_{2}-k)$, if $\lambda_{3N+2}=k$. The other change is when we subtract $k$ from each part,  the partitioned number is  $n-3Nk-3k$ in case of (\ref{4.2}) and the partitioned number is $n-3Nk-2k$ in case of (\ref{4.3}).

The proof of Lemma 4.1 is completed.

Let $|x|<1$, $|y|<1$ and $|q|<1$, we define
\begin{align*}
A_{0}(x,y,q)&=1, \\
A_{N}(x,y,q)&=1+\sum_{\Sigma_{1}=0}^{\infty}\sum_{\Sigma_{2}=0}^{\infty}\sum_{n=1}^{\infty}a_{N}(\Sigma_{1},\Sigma_{2}; n)x^{\Sigma_{1}}y^{\Sigma_{2}}q^{n},\quad N\geq 1.
\end{align*}
The equations in Lemma 4.1 imply the following $q$-difference equations:
\begin{lemma} For all $N\geq 1$
\begin{align}
A_{3N}(x,y,q)=A_{3N-1}(x,y,q)\frac{q^{3N}}{1-q^{3N}}+A_{3N-2}(x,y,q)\frac{q^{3N}}{1-q^{3N}},\\
A_{3N+1}(x,y,q)=A_{3N}(x,y,q)\frac{xq^{3N+1}}{1-xq^{3N+1}}+A_{3N-1}(x,y,q)\frac{xq^{3N+1}}{1-xq^{3N+1}},\\
A_{3N+2}(x,y,q)=A_{3N+1}(x,y,q)\frac{yq^{3N+2}}{1-yq^{3N+2}}+A_{3N}(x,y,q)\frac{yq^{3N+2}}{1-yq^{3N+2}}.
\end{align}
\end{lemma}
Let $|x|<1$, $|y|<1$ and $|q|<1$, if we define
\begin{align*}
P_{0}(x,y,q)&=1, \\
P_{N}(x,y,q)&=1+\sum_{i=1}^{N}\sum_{\Sigma_{1}=1}^{\infty}\sum_{\Sigma_{2}=1}^{\infty}\sum_{n=1}^{\infty}a_{i}(\Sigma_{1},\Sigma_{2};n)x^{\Sigma_{1}}y^{\Sigma_{2}}q^{n}\\
&=1+\sum_{i=1}^{N}A_{i}(x,y;q).
\end{align*}

Thus $P(x,y;q)= \lim_{N\rightarrow \infty} P_{N}(x,y;q) $ is the generating function for the partitions counted by $a(\Sigma_{1}, \Sigma_{2}; n)$. From  Lemma (4.2), we have the following $q$-difference equations for $P_{N}(x,y,q)$.

\begin{theorem} For $N\geq 1$,
\begin{align}
&P_{0}(x,y,q)=1,\nonumber\\
&P_{1}(x,y,q)=\frac{1}{1-xq},\nonumber\\
&P_{2}(x,y,q)=\frac{1}{(1-xq)(1-yq^{2})},\nonumber\\
&P_{3N}(x,y,q)=P_{3N-1}(x,y,q)\frac{1}{1-q^{3N}}-P_{3N-3}(x,y,q)\frac{q^{3N}}{1-q^{3N}},\label{4.7}\\
&P_{3N+1}(x,y,q)=P_{3N}(x,y,q)\frac{1}{1-xq^{3N+1}}-P_{3N-2}(x,y,q)\frac{xq^{3N+1}}{1-xq^{3N+1}},\label{4.8}\\
&P_{3N+2}(x,y,q)=P_{3N+1}(x,y,q)\frac{1}{1-yq^{3N+2}}-P_{3N-1}(x,y,q)\frac{yq^{3N+2}}{1-yq^{3N+2}}.\label{4.9}
\end{align}
\end{theorem}

It is clear that the initial conditions and equations (\ref{4.7}), (\ref{4.8}) and (\ref{4.9}) completely determine $P_{N}(x,y,q)$, hence $P(x,y,q)$.


\noindent {\it Proof.} The initial cases are easy to verify. We prove (\ref{4.7}),  the proofs of (\ref{4.8}) and (\ref{4.9}) are similar.
By using (4.4), 
\begin{align*}
&P_{3N}(x,y,q)\\
&=A_{0}(x,y,q)+A_{1}(x,y,q)+\dots+A_{3N-2}(x,y,q)+A_{3N-1}(x,y,q)+A_{3N}(x,y,q)\\
&=A_{0}(x,y,q)+\dots+A_{3N-2}(x,y,q)+A_{3N-1}(x,y,q)+\left(A_{3N-1}(x,y,q)\frac{q^{3N}}{1-q^{3N}}+A_{3N-2}\frac{q^{3N}}{1-q^{3N}}\right)\\
&=A_{0}(x,y,q)+\dots+A_{3N-3}(x,y,q)+\left(A_{3N-2}(x,y,q)+A_{3N-1}(x,y,q)\right)\frac{1}{1-q^{3N}}\\
&=(A_{0}(x,y,q)+\dots+A_{3N-1}(x,y,q))\frac{1}{1-q^{3N}}-\left(A_{0}(x,y,q) +\dots+A_{3N-3} (x,y,q) \right)\frac{q^{3N}}{1-q^{3N}}\\
&=P_{3N-1}(x,y,q)\frac{1}{1-q^{3N}}-P_{3N-3}(x,y,q)\frac{q^{3N}}{1-q^{3N}}.
\end{align*}

Using $P_{N}(x,y,q)$, we can reformulate our conjecture for moduli three.
\begin{conjecture}
$$
\lim_{N\rightarrow \infty}P_{N}(x,y,q)=\frac{1}{(1-xq)(1-yq^{2})}\lim_{N\rightarrow \infty}P_{N}(x,y,q).
$$
\end{conjecture}



\section{New companions to Rogers-Ramanujan-Andrews-Gordon identities under the conjecture}
Besides Euler's partition theorem involving odd parts and different parts, Rogers-Ramanujan-Andrews-Gordon identities are another famous partition theorem, see \cite{AlladiBerkovich2, AndrewsRR, Ber-peter, Bressoud, Rogers}. Recall that the first Rogers-Ramanujan identity (partition version) says that the number of partitions $n$ with the condition that the difference between any two parts is at least $2$ ( called Rogers-Ramanujan partitions) is equal to the number of partitions of $n$ such that each part congruent to $1$ or $4$ modulo $5$. In view our point,  partitions with each pat congruent to $1$ or $4$ modulo $5$ are exactly partitions belong to $Q$ with length types $(l_{1},0, 0, l_{4})$, $(l_{1},l_{4})\neq (0,0)$. Then our conjecture implies the following companion to the first Rogers-Ramanujan identity:
\begin{theorem}
Suppose the conjecture is true, then the number of partitions of $n$ with with the difference between any two parts is at least
$2$ is equal to the number of partitions of $n$ with parts repeated  at most $4$ times and alternating sum types $(\Sigma_{1},0,0,\Sigma_{4})$,  where $(\Sigma_{1}, \Sigma_{4})\neq (0,0)$.
\end{theorem}
We give an example to illustrate this theorem. Let $n=11$, the partitions of $11$ with the condition that the difference is at least
$2$ are
\begin{align*}
\{11,\quad 10+1,\quad 9+2,\quad 8+3, \quad
7+4, \quad 7+3+1,\quad 6+4+1 \}.
\end{align*}
And the partitions of $11$ with alternating sum types $(\Sigma_{1},0,0,\Sigma_{4})$ are 
\begin{align*}
&\left\{3+2+2+2+2\,(1,0,0,0),\quad 3+2+2+2+1+1\,(2,0,0,1),\quad 4+2+2+2+1\,(2,0,0,1),\right.\\
&\left. 7+1+1+1+1\,(6,0,0,0),\quad    5+2+2+2\,(3,0,0,2),\quad 8+1+1+1\,(7,0,0,1),  \quad 11\,(11,0,0,0)\right\}.
\end{align*}
We list the alternating sum type following each partition.
We have a similar companion on the second Rogers-Ramanujan identity:
\begin{theorem}
Suppose the conjecture is true, then the number of partitions of $n$ with with the difference between any two parts is at least
$2$ and $1$ is not a part is equal to the number of partitions of $n$ with parts repeated  at most $4$ times and alternating sum type $(0,\Sigma_{2},\Sigma_{3},0)$ and $(\Sigma_{2}, \Sigma_{3})\neq (0,0)$.
\end{theorem}
We still use $n=11$ to illustrate this theorem.  The partitions of $11$ with the condition that the difference is at least
$2$ and $1$ is not a part are
\begin{align*}
\{11,\quad 9+2,\quad 8+3,\quad
7+4 \}.
\end{align*}
And the partitions of $11$ with alternating sum types $(0,\Sigma_{2},\Sigma_{3},0)$ are 
\begin{align*}
\left\{3+3+3+1+1\,(0,0,2,0),\, 4+4+1+1+1\,(0,3,0,0),\, 4+4+3\,(0,1,3,0)\, 5+5+1\,(0,4,1,0)\right\}.
\end{align*}
We list the corresponding alternating sum type following each partition.

For Andrews-Gordan's identities, we have 
\begin{theorem}
Let $d\geq 1$, $1\leq i\leq 2d$.
Suppose the conjecture is true, then the number of partitions $\lambda_{1}+\lambda_{2}+\lambda_{3}+\dots+\lambda_{r}$ of $n$ such that no more than $i-1$ of the parts are $1$ and pairs of consecutive integers appear at most $d-1$
 times is equal to the number of partitions of $n$ with parts repeated  at most $2d$ times and alternating sum type $(\Sigma_{1},\Sigma_{2},\dots,\Sigma_{2d-1},\Sigma_{2d})\neq(0,0,\dots,0)$ satisfying that both $\Sigma_{i}$ and $\Sigma_{2d+1-i}$ are zero.
\end{theorem}

\section{Concluding remarks}
In this paper, we give a conjectural generalisation of Euler's partition theorem involving odd parts and different parts for all moduli. We use generating functions and $q$-series analysis to prove the conjecture for two family partitions. Clearly, this method can not be generalised to general cases, it is very difficult to write out the generating functions for partitions in $P$. It is desirable to find combinatorial bijections between two set of partitions  appearing in the conjecture. We note simple combinatorial arguments can not  work in view of our analytic arguments.  Let us consider Theorem 3.2 for $n=14,\, m=3$ and the alternating sum type is $(3, 1)$.  In this case, Theorem 3.2 claims the equation: 
$$  
a(3, 1;14)=a(3,0; 12)+a(3, 0; 9)+a(3, 0; 6)+a(3, 0; 3).
$$
In fact, there are $7$ partitions of $14$ with the alternating sum type $(3,1)$  and they and their conjugates are
\begin{align*}
&7+4+3 &\longleftrightarrow \quad\quad\quad\quad&3+3+3+{\color{blue} 2}+1+1+1\\
&4+4+3+3+0+0 &\longleftrightarrow\quad\quad\quad\quad &4+4+4+{\color{blue} 2}\\
&5+4+3+2+0+0 &\longleftrightarrow\quad\quad\quad\quad &4+4+3+{\color{blue} 2}+1\\
&6+4+3+1+0+0 &\longleftrightarrow\quad\quad\quad\quad&4+3+3+{\color{blue} 2}+1+1\\
&5+3+3+2+1+0 &\longleftrightarrow\quad\quad\quad\quad &{\color{blue} 5}+4+3+1+1\\
&6+3+3+1+1+0 &\longleftrightarrow\quad\quad\quad\quad &{\color{blue} 5}+3+3+1+1+1\\
&5+3+2+2+1+1 &\longleftrightarrow\quad\quad\quad\quad &6+4+{\color{blue} 2}+1+1.
\end{align*}
There are $3$ partitions of $12$ with the alternating sum type $(3,0)$  and their conjugates are
\begin{align*}
&6+3+3 \quad\quad\quad&\longleftrightarrow\quad\quad\quad\quad &3+3+3+1+1+1\\
&4+3+3+2+0+0 &\longleftrightarrow\quad\quad\quad\quad &4+4+3+1\\
&5+3+3+1+0+0 &\longleftrightarrow\quad\quad\quad\quad &4+3+3+1+1.
\end{align*}
There are $2,1,1$ partitions of $9,6,3$ respectively with the alternating sum type $(3,0)$  and their conjugates are
\begin{align*}
5+2+&2 &\longleftrightarrow& \quad\quad\quad&3&+3+1+1+1,\\
4+2+&2+1+0+0 &\longleftrightarrow&\quad\quad\quad\quad &4&+3+1+1,\,\quad\quad\\
4+1+&1&\longleftrightarrow&\quad\quad\quad\quad&3&+1+1+1,\,\quad\quad\\
3+0+&0&\longleftrightarrow&\quad\quad\quad\quad&1&+1+1.\,\,\,\,\,\quad\quad\quad
\end{align*}
It seems that each partition of $14$  with the alternating sum type $(3,1)$ can be obtained from the partitions of $12,9,6,3$
with the alternating sum type $(3,0)$ through the following steps:

$\bullet$ \textbf{Step 1.} Change each partition $\lambda$ of $12, 9,6,3$ into its conjugate $\mu$,

$\bullet$ \textbf{Step 2.}  Add 2 to $\mu$, as a new part of $\mu$, if $|\lambda|=12$;  add 5 to $\mu$ as a new part of $\mu$, if $|\lambda|=9$; add 8 to $\mu$ as a new part, if $|\lambda|=6$; and add $11$ to $\mu$ as a new part, if $|\lambda|=3$.  Denote $\mu^{\prime}$ for the obtained partition in each case.


$\bullet$ \textbf{Step 3.} Change each $\mu^{\prime}$ into its conjugate.

However,  we have two questions arising from this combinatorial arguments for general $n$.  One is some partitions of $n$ can not be obtained in this way.  For example,  if you omit the part $2$ from the partition $4+4+4+2$,  you get the partition $4+4+4$, which does not belong to  conjugates of any partitions of $12$ with the alternating sum type $(3,0)$, since  each conjugate of a partition of $12$ with type $(3, 0 )$ has at least $3$ as its part !  Hence $4+4+4+2$ can not be obtained by the steps above.
Second question is in \textbf{Step 3}, in some cases,  the conjugate of $\mu^{\prime}$ may have parts repeated more than 3 times! 





\section*{Acknowledgements}
The author would like to thank Professor Peter Paule  for their valuable comments on an earlier version of this paper and encouragements. 

\begin{thebibliography}{99}
\bibitem{Alder} H.L. Alder, \textit{Partition identities - from Euler to the present}, Am. Math. Monthly  \textbf{76} (1969),733-746.

\bibitem{AlladiWeighted} K. Alladi, \textit{Partition identities involving gaps and weights}, Trans. Amer. Math. Soc. \textbf{349} (1997), no. 12, 5001-5019.

\bibitem{All3} K. Alladi, \textit{Partitions with non-repeating odd parts and combinatorial identities}, Ann. Comb. \textbf{20} (2016), 1-20.

\bibitem{AlladiBerkovich2} K. Alladi, and A. Berkovich, \textit{New weighted Rogers-Ramanujan partition theorems and their implications}, Trans. Amer. Math. Soc. \textbf{354} (2002), no. 7, 2557-2577.

\bibitem{AndrewsRR} G. E. Andrews, \textit{An analytical proof of the Rogers-Ramanujan-Gordon identities}, Am. J. Math. \textbf{88} (1966),  844-846.    

\bibitem{Theory_of_Partitions} G. E. Andrews, \textit{The theory of partitions}, Cambridge Mathematical Library, Cambridge University Press, Cambridge, 1998. 

\bibitem{AndrewsEuler} G. E. Andrews, \textit{On generalizations of Euler's partition theorem}, Michigan Math. J. \textbf{1}
 (1966), 491--498.
\bibitem{LostNotebook_1} G. E. Andrews, and B. C. Berndt \textit{Ramanujan's lost notebook. Part I}, Springer, New York, 2005. 

\bibitem{LostNotebook_2} G. E. Andrews, and B. C. Berndt \textit{Ramanujan's lost notebook. Part II}, Springer, New York, 2008. 


\bibitem{Ber-peter}
A. Berkovich and P. Paule, \textit{Variants of the Andrews-Gordon identities}. Ramanujan J. \textbf{5} (2001), no. 4, 391-404.

\bibitem{bes91} C. Bessenrodt, \textit{A combinabinatorial proof of a refinement of
the Andrews-Olsson partition identity}, European J. Combin.  \textbf{12} (1991) 271--276.

\bibitem{bes94} C. Bessenrodt, \textit{A bijection for Lebesgue's partition
identity in the spirit of Sylvester}, Discrete Math.  \textbf{132} (1994) 1--10.

\bibitem{bou97}M. Bousquet-M\'{e}lou and K. Eriksson,
\textit{Lecture hall partitions}, Ramanujan J.  1 (1997) 101--111.

\bibitem{bou972}M. Bousquet-M\'{e}lou and K. Eriksson,
\textit{Lecture hall partitions II}, Ramanujan J.  1 (1997) 165--185.



\bibitem{Bressoud} D.M. Bressoud, \textit{A generalization of the Rogers-Ramanujan identities for all moduli}, J. Combin. Theory Ser. A  \textbf{27} (1979), no. 1, 64-68. 


\bibitem{Chen} W.Y.C. Chen, H.Y. Gao, K.Q. Ji, M.Y.X. Li, \textit{A unification of two refinements of Euler's partition theorem}, Ramanujan J. \textbf{23} (2010), 137-149.


\bibitem{Fine} N. J. Fine, \textit{Some new results on partitions}, PNAS,  \textbf{34} (1948),616-618.

\bibitem{GasperRahman} G. Gasper, and M. Rahman, \textit{Basic Hypergeometric Series}, Vol. 96. Cambridge university press, 2004.

\bibitem{Gollnitz} H. G\"ollnitz, \textit{Partitionen mit Differenzenbedingungen}, J. reine angew. Math. \textbf{225} (1967), 154-190.

\bibitem{Gordon} B. Gordon, \textit{A combinatorial generalization of the Rogers-Ramanujan identities}, Duke Math. J. \textbf{32} (1965), 741-748.


\bibitem{Kim-Yee} D. Kim and A. J. Yee, \textit{A note on partitions into
distinct parts and odd parts}, Ramanujan J. \textbf{3} (1999) 227--231.

\bibitem{Pak} I. Pak, \textit{Partition bijections, a survey}, Ramanujan J. \textbf{12} (2006), 5-75.

\bibitem{Pak-Postnikov} I. Pak and A. Postnikov, \textit{A generalization of Sylvester's Identity}, Discrete Math. \textbf{178} (1998), 277-281.


\bibitem{Rogers} L. J. Rogers, \textit{On two theorems of combinatory analysis and some allied identities}, Proc. London Math. Soc. \textbf{16} (1917), 315-336.

\bibitem{SavageYee} C.D. Savage, A.J. Yee, \textit{Euler's partition theorem and the combinatorics of $l$-sequences}. J. Combin. Theory Ser. A \textbf{115} (2008), 967-996.



\bibitem{Syl82} J. Sylvester,  \textit{A constructive theory of partitions, arranged in three acts, an interact and an exodion}, Amer. J. Math. \textbf{5} (1882), 251--330.   


\bibitem{Yee2001} A.J.Yee, \textit{On the combinatorics of lecture hall partitions}, Ramanujan J. \textbf{5} (2001), 247-262.

\bibitem{Yee2002} A.J. Yee, \textit{On the refined lecture hall theorem}, Discrete Math. \textbf{248} (2002),  293-298.

\bibitem{Zeng} Jiang Zeng, \textit{The q-variations of Sylvester's bijection between odd and strict partitions}, Ramanujan J. \textbf{9} (2005), 289-303.
\end{thebibliography}



\end{document}

